\documentclass[11pt]{article}
\usepackage[margin=1.1in]{geometry}
\usepackage{amsmath,amssymb,amsthm}
\usepackage{xcolor}
\usepackage[T1,OT1]{fontenc}
\usepackage[var0,noupquote]{zi4}\renewcommand{\ttdefault}{cmtt}
\newcommand{\bin}[1]{{\fontencoding{T1}\fontfamily{zi4}\selectfont #1}}% binary strings: a monospace whose bold stands out
\usepackage{longtable,array}
\usepackage[colorlinks=true,linkcolor=blue,citecolor=blue,urlcolor=blue]{hyperref}

\newtheorem{theorem}{Theorem}[section]
\newtheorem{proposition}[theorem]{Proposition}
\newtheorem{lemma}[theorem]{Lemma}
\newtheorem{corollary}[theorem]{Corollary}
\newtheorem{definition}[theorem]{Definition}
\newtheorem{hypothesis}[theorem]{Hypothesis}
\newtheorem{heuristic}[theorem]{Heuristic}
\theoremstyle{remark}
\newtheorem{remark}[theorem]{Remark}

\newcommand{\Z}{\mathbb{Z}}
\newcommand{\vt}[1]{v_2\!\left(#1\right)}
\newcommand{\vth}[1]{v_3\!\left(#1\right)}
\newcommand{\w}{\omega}
\newcommand{\wnext}{\w_{+}}
\newcommand{\dnext}{d_{+}}
\newcommand{\LL}{\log_2 3}
\newcommand{\kp}{\kappa_p}
\newcommand{\stot}{s_{\mathrm{tot}}}
\newcommand{\Ent}{\mathcal{H}}
\setlength{\emergencystretch}{3em}
% DOI: replace the placeholder below with the v4 version-specific DOI once it is minted
% (it is used on the title page and in the Version note).
\newcommand{\vfourdoi}{\textbf{[v4 version DOI]}}
% PIN-AT-BUILD: replace with the commit of the public record at which every wiki section and
% script named in this paper is cited (the record's main is not yet pushed; do not cite an
% unreachable commit as reachable). Used in Appendix D ("Verification and record") and Appendix A.
\newcommand{\recordcommit}{\textbf{[record commit, pinned at release]}}

\title{Reduced coordinates for the Collatz map:\\ exact per-step laws, anchor dynamics,\\ and the limits of counting arguments for cycles}
\author{Ben Macindoe\thanks{\texttt{macindoebenjamin@gmail.com}. Developed through extensive AI-assisted research, with ChatGPT (OpenAI) from December 2025 to March 2026 and with Claude (Anthropic) from March 2026; the division of labor and the verification protocol are stated in Appendix~\ref{app:method}.}}
% PIN-AT-BUILD: the month is the month of minting.
\date{v4, October 2026 $\;\cdot\;$ DOI: \vfourdoi{} $\;\cdot\;$ concept DOI: \href{https://doi.org/10.5281/zenodo.21273547}{10.5281/zenodo.21273547}}

\hypersetup{
  pdftitle={Reduced coordinates for the Collatz map: exact per-step laws, anchor dynamics, and the limits of counting arguments for cycles},
  pdfauthor={Ben Macindoe}
}

\begin{document}
\maketitle

\begin{abstract}
The Collatz conjecture says that if you repeatedly halve an even number, or multiply an odd number by $3$ and add $1$, you always reach $1$. Much of that process is forced bookkeeping: as long as one more than the current odd number is divisible by $4$, each step simply trades a factor of $2$ for a factor of $3$. This paper folds each such forced run into a single \emph{block}, and describes each block by two whole numbers, its \emph{state}. (The same states were found independently by Williams, 2026.) The states step from one to the next by a new map. That map loses nothing: it gives the same answers as the original rule to whether every orbit (the sequence of numbers a start visits) reaches $1$ and to which loops (cycles) exist.

For one step of the new map the paper proves exact laws for four things: how many halvings end a block; when a factor of $3$ is picked up; how deep the next block runs; and how a single fixed binary string attached to each state, its \emph{anchor}, changes from one state to the next. All of these except the factor of $3$ are governed by the anchor; the factor of $3$ obeys its own exact law. A finite number of binary digits of the state, with three numbers describing the step (Section~\ref{sec:anchor}), decides each step: it either gives the answer or says correctly that the answer lies deeper, and it is never wrong. An accounting, stated as a heuristic rather than proved, indicates that no fixed finite number of digits can decide an orbit's whole future, which locates the difficulty in where the anchors' digits come from.

On cycles, a one-line identity rewrites a cycle as whole-number conditions of two kinds: \emph{size} conditions (about how close a power of $2$ must come to a power of $3$) and \emph{divisibility} conditions. From it come short new derivations of two classical facts: there are no cycles other than $1, 4, 2, 1$ made of one, two or three blocks; and the classical bound of Rozier (1991) and Simons--de~Weger (2005) on how many odd steps (multiplications by $3$) a cycle with a given number of blocks can have, whose strength fades exponentially as the number of blocks grows. The paper proves that this fading is unavoidable for any argument from size alone. An explicit family of near-cycles, the \emph{staircases} (profiles of an ever-growing orbit bent into a loop), passes every size test, for every choice of starting block, while its measure of closeness (how small the gap between the two powers is, compared with the larger one) stays between two fixed numbers: at every number of blocks from $3$ on (from $16$ on by proof, $3$ to $15$ by computer check; for $2$ and $4$ blocks passing examples are shown, with no such range claimed). When every number in the cycle must be at least a given size, the staircases pass at every number of blocks from an explicit threshold on, with their measure of closeness exceeding the classical bound's own exponent by at most a fixed amount. At that threshold scale they attain the optimum of Rozier's 1991 analysis --- the best worst case over block depths allowed to be any real numbers --- under every size test, the sharpness Simons and de~Weger expected in 2005. The analogue for cycle length, and the conclusion that ruling out cycles of every length at once needs divisibility and not size alone, are Halbeisen and Hungerb\"uhler's (1997).

Finally the paper states the `fair-coin' assumption behind the classical heuristics as a precise conjecture. For large starting values chosen at random, record each block's (entry depth, exit valuation) pair while the orbit has used up fewer halvings than a fixed multiple of the start's number of binary digits, over a number of blocks that is itself a fixed multiple of that number of digits. The conjecture is that these pairs occur, for patterns of every finite length, with the frequencies of independent fair-coin counts --- for every admissible choice of the two multiples (one that fits the record inside the halving budget on the model's average and keeps all but a vanishing share of orbits away from small numbers; Section~\ref{sec:aeh}). It proves what follows if the conjecture holds, shows that the fair-coin conjecture holds outright while the orbit is still within the start's own supply of binary digits (the proof is assembled in the project record), and reports a test campaign whose four apparent anomalies all dissolved under sharper controls. The technical summary at the start of the appendices states all of this in the usual notation.
\end{abstract}


\subsection*{Version note}
% DOI: the version-specific DOI enters through \vfourdoi (preamble).
% PIN-AT-BUILD: the month below is the month of minting.
v4, October 2026 (DOI \vfourdoi). It supersedes v3 (DOI 10.5281/zenodo.21730505, August 2026), v2 (10.5281/zenodo.21421120) and v1 (10.5281/zenodo.21273548, July 2026); the concept DOI 10.5281/zenodo.21273547 resolves to the latest version. v4 attributes the uniform trim (Theorem~\ref{thm:uniform}) to Rozier and to Simons--de~Weger, with finiteness at each period to Rozier and Brox, and presents it as a rederivation; states the sharpness of that trim (Theorem~\ref{thm:staircase}) at every number of blocks from $3$ on, and at every size threshold from an explicit threshold number of blocks on, with its proof, in place of the assessment and status report of v3; cites the classical and contemporary sources of the reduced coordinates' objects and of the elimination identity; and corrects the attribution of Theorem~\ref{thm:vlaw}'s unconditional bound at the odd core $1$. v4 is also rewritten for a curious reader without training in mathematics: the main text carries the ideas in words, and the statements, with their proofs or the sections of the project record that hold them, are in Appendices~\ref{app:reduced}--\ref{app:aeh}. Its statements are therefore renumbered, and Appendix~\ref{app:concordance} maps v3's numbers to v4's. For outside citations of this paper's statements by number, v3 remains the pinned citation. The history of every version, item by item, is at \texttt{paper/collatz-reduced-version-history.md} in the project record.

\subsection*{Author's note}

\emph{Dear reader --- I'm no mathematician, and I learned much of this paper's terminology during the process of formalizing it ($p$-adic valuations were completely new to me as a formal concept). It all began when I got hooked on a strange recursive pattern in the prime factors of a Collatz chain, and that idea pushed me into conversation with AI.
We cleaned my idea up into the states $(\w,d)$, with the primary goal of factoring as much deterministic behaviour as possible out of the question, in the hope of focusing attention where it matters.
What these coordinates contribute is modest: an exact description of each step, a construction agreeing that Rozier and Simons--de~Weger's limit on counting arguments against cycles cannot be substantially improved, and a precise statement of the `fair-coin' assumption behind the classical heuristics. Because this work is AI-assisted and because I am not formally trained in the relevant mathematics, Appendix~\ref{app:method} describes the verification protocol used throughout the project. The complete research record is public at \url{https://github.com/macindoe/collatz}. --- Ben}

% =====================================================================
%  MAIN TEXT
% =====================================================================

\section*{Symbols at a glance}
\addcontentsline{toc}{section}{Symbols at a glance}

This page lists every symbol used outside the appendices --- in the front matter and in Sections~\ref{sec:intro}--\ref{sec:credit} --- with its meaning in words and the section that introduces it. Each is also explained in words where it first appears. Section~\ref{sec:intro} and the abstract use none of them. (The letter $p$ also appears in the fixed phrase `$p$-adic', in the Author's note. The appendices use many more symbols, among them $q$, $\gamma$, $X_0$ and $p_1(X_0)$ of Appendix~\ref{app:cycles}'s guided account.) Every symbol of the paper, including those of the appendices, is in Appendix~\ref{app:symbols}.

\begin{center}
\begin{tabular}{>{\raggedright\arraybackslash}p{0.12\textwidth}>{\raggedright\arraybackslash}p{0.6\textwidth}>{\raggedright\arraybackslash}p{0.18\textwidth}}
\hline
Symbol & Meaning & Introduced \\
\hline
$x$ & an odd number of an orbit; in the running example, $x = 7$ & \S\ref{sec:block} \\
$T$ & the Collatz map on odd numbers: multiply by $3$, add $1$, halve until odd & \S\ref{sec:block}; App.~\ref{app:reduced} \\
$m$ & a block's \emph{entry depth}: how many times $2$ divides $x+1$ (the number of trailing $1$s of $x$ in binary) & \S\ref{sec:block} \\
$s$ & the \emph{exit valuation}: how many halvings, beyond the first, end every block of a given state & \S\ref{sec:block} \\
$a$ & how many factors of $3$ are left in $x+1$ after the factors of $2$ are removed & \S\ref{sec:reduced} \\
$\w$ & the \emph{odd core}: what is left of $x+1$ after removing every factor of $2$ and of $3$ & \S\ref{sec:reduced} \\
$d$ & the \emph{depth}: $m$ plus $a$ & \S\ref{sec:reduced} \\
$(\w,d)$ & a block's \emph{state} & \S\ref{sec:reduced}; Def.~\ref{def:reduced} \\
$F$ & the \emph{reduced map}, from each state to the next & \S\ref{sec:reduced}; Def.~\ref{def:reduced} \\
$p$ & the number of blocks in a cycle & \S\ref{sec:cycles} \\
$n$ & the sum of a cycle's entry depths: how many times it multiplies by $3$ & \S\ref{sec:cycles} \\
$K$ & a cycle's total number of halvings: $n$ plus the sum of its exit valuations & \S\ref{sec:cycles} \\
$\gamma$ & a cycle's closeness measure: the gap between $2^K$ and $3^n$ is $2^K$ divided by $2^{\gamma}$ & \S\ref{sec:cycles} \\
$m|s$ & a \emph{letter}: a block's entry depth and exit valuation, written with a bar so as not to be mistaken for a state & \S\ref{sec:aeh} \\
$\lambda$, $k$ & Williams's names for $\w$ and $d$ (Section~\ref{sec:credit} only) & \S\ref{sec:credit} \\
\hline
\end{tabular}
\end{center}

Three words are used as names: the \emph{anchor} of $\w$, a fixed binary string, running on forever to the left, that belongs to $\w$ (all zeros when $\w = 1$; Definition~\ref{def:anchor} defines it); a \emph{lifting} state, one whose $\w$ leaves remainder $1$ on division by $8$ with $d$ even, or remainder $3$ with $d$ odd; and the \emph{absorption}, the number of factors of $3$ the next state picks up (all in Section~\ref{sec:anchor}).

\section{What this paper is about}\label{sec:intro}

Take a positive whole number. If it is even, halve it; if it is odd, multiply it by $3$ and add $1$. Repeat. The Collatz conjecture says that from every starting number this eventually reaches $1$, after which it goes round $1, 4, 2, 1$ forever. Two things could go wrong. An \emph{orbit} (the sequence of numbers visited) could grow forever; or it could fall into a different \emph{cycle} --- a loop of numbers repeating forever without reaching $1$.

This paper has three purposes. The first is to prove exact laws for each step of the process, once its forced bookkeeping has been folded away. The second is to determine how far counting arguments --- arguments from how large the numbers in a cycle would have to be --- can go toward ruling cycles out. The classical answer: they bound how many odd steps a cycle can have, but the bound weakens exponentially as the number of \emph{blocks} (forced runs of steps, defined in Section~\ref{sec:block}) in the cycle grows \cite{rozier91,sdw}. The paper proves that size arguments go no farther: this exponential weakening cannot be avoided, for every number of blocks from $3$ on. It also proves the same when every number in the cycle must be at least a given size, for every number of blocks from an explicit threshold on (the threshold depends on the size). The third purpose is to state precisely the `fair-coin' assumption --- that the halvings behave like tosses of a fair coin --- which the classical heuristic arguments leave implicit.

\paragraph{Three levels of description.} The paper looks at the process at three levels, one at a time.
\begin{enumerate}
\item[(1)] \emph{The map, one step at a time.} After an odd step the result is always even ($3$ times an odd number is odd, and an odd number plus $1$ is even), so at least one halving always follows. Because of that, the paper does all the halvings at once and looks only at odd numbers: one step multiplies by $3$, adds $1$, and halves until the result is odd again. (This is Tao's `Syracuse map' \cite{tao}.)
\item[(2)] \emph{Blocks.} Many steps are forced: while one more than the number is divisible by $4$, each step trades a $2$ for a $3$; otherwise one burst of halvings ends the run. Section~\ref{sec:block} shows how on an example, and calls the whole run a block.
\item[(3)] \emph{States.} Each block is described by two whole numbers, its \emph{state} (Section~\ref{sec:reduced}). The blocks chain into a map from each state to the next, the \emph{reduced map}. It loses nothing the conjecture depends on: its orbits reach the state of the number $1$ exactly when the original orbits reach $1$; the trivial cycle $1, 4, 2, 1$ corresponds to the state of $1$, which the reduced map sends to itself; and every other cycle of numbers corresponds to exactly one other cycle of states, and back.
\end{enumerate}
The rest of the paper works at the third level: exact laws for one step (Section~\ref{sec:anchor}); what counting can and cannot say about cycles (Section~\ref{sec:cycles}); and the fair-coin assumption (Section~\ref{sec:aeh}). Section~\ref{sec:open} says what remains open, and Section~\ref{sec:credit} who did what. Section~\ref{sec:block}'s example, the orbit of $7$, runs through Sections~\ref{sec:block}--\ref{sec:anchor}, and its blocks return in Section~\ref{sec:aeh}.

\paragraph{What is new.} Each item is explained in the section named, and stated formally in the appendix named.
\begin{enumerate}
\item[(i)] \emph{Compression without loss.} Each forced run of steps becomes a single step of the reduced map. The conjecture's two questions --- does every orbit reach $1$, and is there any cycle other than $1, 4, 2, 1$ --- have exactly the same answers for the reduced map as for the original one. The paper proves that the reduced map is well defined and that this equivalence holds. (The states were found independently by Williams \cite{williams}, who calls them skeleton rows. Section~\ref{sec:reduced}; Appendix~\ref{app:reduced}.)
\item[(ii)] \emph{How far each block falls.} One exact rule gives the number of halvings, beyond the first, that end each block. Section~\ref{sec:anchor} sorts the states into eight kinds by two simple remainders. For six of the kinds the number is fixed. For the other two it is read off by comparing the state's depth, digit by digit, with a fixed binary pattern belonging to the state. (This is the exit rule of Section~\ref{sec:anchor}; Theorem~\ref{thm:vlaw}.)
\item[(iii)] \emph{When factors of $3$ appear, and how deep the next block runs.} Exact formulas, valid for every state, say when the next state gains a factor of $3$ and how deep the next block is. (Lemma~\ref{lem:absorption}, Theorem~\ref{thm:depth}.)
\item[(iv)] \emph{How much must be known to predict a step.} Finitely many binary digits of the state, with three numbers describing the step (how far it falls, how deep the next block enters, and how many factors of $3$ it picks up), either predict the next step exactly or say, correctly, that the next burst of halvings is longer than those digits can determine. They never give a wrong answer. How fast those digits are used up is estimated by an accounting stated as a heuristic --- a reasoned estimate, not a proof. (Theorems~\ref{thm:deltaM}--\ref{thm:onestep}, Heuristic~\ref{prop:budget}.)
\item[(v)] \emph{What counting can and cannot do against cycles.} A one-line identity (stated in Section~\ref{sec:cycles}) turns a cycle into conditions of two kinds: size (one number must be at most another) and divisibility (one must divide another). From it come short new derivations of classical results: no cycles other than $1, 4, 2, 1$ with one, two or three blocks; and the uniform bound of Rozier and Simons--de~Weger, whose strength fades exponentially as the number of blocks grows. The paper then proves that this fading is unavoidable for size arguments. An explicit family of near-cycles, the \emph{staircases}, meets every size condition, for every choice of starting block, at every number of blocks from $3$ on (by proof from $16$ on, by computer check for $3$ to $15$; for $2$ and $4$ blocks, examples only). When every number must be at least a given size, it meets them at every number of blocks from an explicit threshold on, and there it reaches the optimum of Rozier's analysis --- the best worst case over block depths allowed to be any real numbers. (Section~\ref{sec:cycles}; Appendix~\ref{app:cycles}.)
\item[(vi)] \emph{The fair-coin assumption made precise.} The statistical assumption behind the classical heuristics is stated exactly. Its consequences are proved on the condition that it holds, it is proved outright in the range where a start's own binary digits still decide its blocks, and the tests run on it are reported. (Section~\ref{sec:aeh}; Appendix~\ref{app:aeh}.)
\end{enumerate}

\paragraph{What this paper does not claim.}
\begin{itemize}
\item It does not settle the Collatz conjecture.
\item It does not prove that no fixed, finite number of binary digits can decide an orbit's whole future (Section~\ref{sec:anchor}). That is the paper's organizing heuristic: a guiding estimate, not a proved theorem.
\item The exclusion of cycles with one, two or three blocks is classical (Steiner; Rozier 1990; Simons--de~Weger; Hercher). So is the uniform bound on cycles that Section~\ref{sec:cycles} calls the \emph{trim} (Rozier 1991; Simons--de~Weger); and Rozier (1991) and Brox (2000) proved that only finitely many cycles can have any given number of blocks. What is offered for them is new derivations.
\item Every statement about the staircases concerns how far size arguments reach, and none bears on whether cycles exist. Nothing is stated for one block. For two and four blocks, only examples that pass the size tests are shown, with no claim about their closeness measure. The threshold number of blocks is computed for each size, but no single formula covering every size is claimed, and nothing is claimed below the threshold.
\item The fair-coin hypothesis is a conjecture, apart from the range where it is proved (Section~\ref{sec:aeh}). Even where it holds, no statement that orbits drift downward or shrink follows from it. The paper does not state, as a consequence of it, that almost every orbit falls below its start (`descent'): that is already proved without it, and more strongly, by Inselmann \cite{inselmann}. The starting values whose blocks it allows to stray from the coin pattern form a fraction of all starting values that shrinks to zero as the starting size grows; that is not the same as saying that almost no orbits stray. It says nothing about where any single orbit ends up in the long run, does not rule out a single orbit shaped like a staircase, and cannot be applied again and again along one orbit.
\end{itemize}

\paragraph{How to read this paper.} Sections~\ref{sec:block} and~\ref{sec:reduced} build the coordinates on one worked example. Sections~\ref{sec:anchor}, \ref{sec:cycles} and~\ref{sec:aeh} are one page each: what is claimed, why it matters, one worked example, and where the full account lives. Each of Appendices~\ref{app:laws}, \ref{app:cycles} and~\ref{app:aeh} opens with a guided account in plain words before its formal statements. ``Symbols at a glance'', before this section, lists every symbol the main text uses. The main text uses only short formulas, each explained in words; the longer formulas are in the appendices. Each result stated informally in the main text names its formal statement in an appendix, where definitions, constants, proofs and pointers into the project record are given in full: Appendix~\ref{app:reduced} (the reduced system), \ref{app:laws} (the exact laws), \ref{app:cycles} (cycles) and \ref{app:aeh} (the hypothesis). The technical summary, at the start of the appendices, is the abstract in the usual notation. Appendix~\ref{app:method} states how the work was done and checked; Appendix~\ref{app:related} discusses related work in detail; Appendix~\ref{app:verify} indexes every computational claim to the script that checks it; Appendix~\ref{app:concordance} maps v3's statement numbers to this version's; and Appendix~\ref{app:symbols} indexes every symbol of the paper.

\section{One rising run: trading 2s for 3s}\label{sec:block}

\emph{What this section is about.} Long stretches of a Collatz orbit are forced. This section shows the mechanism on one example, then states it in general. Everything else in the paper is built on it.

From here on the paper writes $T$ for the map on odd numbers of Section~\ref{sec:intro}: multiply by $3$, add $1$, and halve until odd. Start from $x = 7$ and write each odd number of its orbit, and the number one above it, in binary (the cascade row is the even number $26$, half of $3 \cdot 17 + 1$):
\begin{center}
\begin{tabular}{rrrrl}
\hline
step & $x$ & $x$ in binary & $x+1$ in binary & $x + 1$ factored \\
\hline
0 & 7 & \bin{\textbf{111}} & \bin{1\textbf{000}} & $2^3\cdot 3^0\cdot 1$ \\
1 & 11 & \bin{10\textbf{11}} & \bin{11\textbf{00}} & $2^2\cdot 3^1\cdot 1$ \\
2 & 17 & \bin{1000\textbf{1}} & \bin{1001\textbf{0}} & $2^1\cdot 3^2\cdot 1$ \\
\hline
cascade & 26 & \bin{11010} & \bin{11011} & $2^0\cdot 3^3\cdot 1$ \\
\hline
3 & 13 & \bin{110\textbf{1}} & \bin{111\textbf{0}} & $2^1\cdot 3^0\cdot 7$ \\
\hline
\end{tabular}
\end{center}
The steps from $7$ to $11$ and from $11$ to $17$ are \emph{rising} steps: $3x+1$ can be halved only once, and half of $3x+1$ is more than $x$, so the result is larger than where it started. Here is why it halves only once. When $x+1$ is divisible by $4$, so is $3$ times $x+1$; and $3x+1$ is that number minus $2$, so it is $2$ more than a multiple of $4$, and can be halved only once (a number $2$ more than a multiple of $4$ is twice an odd number). (For $x = 7$: $3 \cdot 8 = 24$ is a multiple of $4$, and $22 = 24 - 2$ halves once, to $11$.)

A rising step trades a factor of $2$ for a factor of $3$. After the step, the new number plus one is exactly three halves of the old number plus one: $11 + 1 = 12$ is three halves of $7 + 1 = 8$, and $17 + 1 = 18$ is three halves of $12$. So each rising step removes one factor of $2$ from $x+1$ and puts in one factor of $3$, as the last column shows. In binary, $x+1$ loses one trailing zero per step (\bin{1\textbf{000}}, \bin{11\textbf{00}}, \bin{1001\textbf{0}}), and so $x$, which is one less, loses one trailing $1$ (\bin{\textbf{111}}, \bin{10\textbf{11}}, \bin{1000\textbf{1}}).

The trade stops when only one factor of $2$ is left in $x+1$, as at $17$, where $x + 1 = 18$. Then $3$ times $x+1$ is twice an odd number, and subtracting $2$ leaves a multiple of $4$, so the step halves at least twice: $3 \cdot 17 + 1 = 52$, which is \bin{110100} in binary, halves twice, down to $13$.

In general (Proposition~\ref{prop:block}): let $m$ be the number of times $2$ divides $x+1$ --- equivalently, the number of trailing $1$s of $x$ in binary. The next $m - 1$ steps are rising steps, each converting one factor of $2$ in $x+1$ into a factor of $3$. The step after them ends in a \emph{cascade}: a burst of at least two halvings, down to an odd number, the block's \emph{exit}. Write $s$ for the number of the cascade's halvings, not counting its first; the paper calls it the block's \emph{exit valuation}. The whole run --- $m$ odd numbers, $m-1$ rising steps and one cascade --- is a \emph{block}, and $m$ is its \emph{entry depth}. For $x = 7$: $m = 3$, the block is $7, 11, 17$, and the cascade $52 \to 26 \to 13$ halves twice, so $s = 1$ and the exit is $13$. (Appendix~\ref{app:reduced} gives the exit by formula.)

This mechanism is classical (Steiner \cite{steiner}; see also Rozier \cite{rozier90} and LaDue \cite{ladue}), and it is also the mechanism of the author's hand-drawn note of March~2023 (Section~\ref{sec:credit}).

\emph{The same block in the one-division map.} Much of the literature uses a slower map, the \emph{one-division map} (Terras \cite{terras}): from an odd number, multiply by $3$, add $1$ and halve once; from an even number, halve once. In it, a block of entry depth $m$ and exit valuation $s$ is exactly $m$ odd numbers followed by $s$ even ones: for $x = 7$, the odd numbers $7, 11, 17$, then the even number $26$, then the next block starts at $13$. So results stated for the one-division map can be read in blocks (Appendix~\ref{app:cycles}, the dictionary paragraph).

\section{Blocks become states}\label{sec:reduced}

\emph{What this section is about.} A block is determined by fewer numbers than it seems: two, its state. This section explains why, and states the theorem that nothing the conjecture depends on is lost. The formal versions are in Appendix~\ref{app:reduced}.

\paragraph{The state.} Remove the $m$ factors of $2$ from $x+1$. What is left is odd; remove from it every factor of $3$ as well, and write $a$ for how many there were. What is left now, $\w$, the \emph{odd core}, is divisible by neither $2$ nor $3$. The \emph{depth} $d$ is $m$ plus $a$, and the pair $(\w, d)$ is the block's \emph{state}. For $x = 7$: $x + 1 = 8$ has $m = 3$, and only $1$ is left, so $a = 0$, $\w = 1$, $d = 3$, and the state is $(1,3)$.

\paragraph{Why the state is enough.} After the $m-1$ rising steps, one more than the current number is $2$ times $\w$ times $3$ to the power $d - 1$. The reason: $x+1$ had $a$ factors of $3$, and each of the $m - 1$ rising steps added one more, $d-1$ in all. That number depends only on $\w$ and $d$, so the cascade, and the exit, depend only on them too (Proposition~\ref{prop:block}(ii)). In the example, $17 + 1 = 18$ is $2$ times $\w = 1$ times $3^2$, with $d - 1 = 2$ threes.

So different odd numbers can share a state and exit at the same place; call them the state's \emph{representatives}. Exactly $d$ odd numbers have a given state $(\w, d)$, one for each way of writing $d$ as $m$ plus $a$ with $m$ at least $1$ (Proposition~\ref{prop:block}(iii)). For $(1,3)$ they are $7$ ($m = 3$), $11$ ($m = 2$, as $12$ is $4 \cdot 3$) and $17$ ($m = 1$, as $18$ is $2 \cdot 9$): the block of Section~\ref{sec:block} consists of exactly the numbers of its own state, and all three exit at $13$: $7$, $11$ and $17$ all end in the same cascade.

\paragraph{The reduced map.} The exit is again odd and begins the next block, whose state is found the same way. The \emph{reduced map} $F$ sends each state to the state of the next block; Definition~\ref{def:reduced} gives it by formulas. For $x = 7$:
\begin{itemize}
\item the block $7, 11, 17$ has state $(1,3)$ and exits at $13$;
\item $13 + 1 = 14$ is $2 \cdot 7$, so the next state is $(7,1)$: a block of one odd number ($m = 1$), and $3 \cdot 13 + 1 = 40$ halves three times ($s = 2$), to $5$;
\item $5 + 1 = 6$ is $2 \cdot 3$, so the next state is $(1,2)$: again a block of one odd number, and $3 \cdot 5 + 1 = 16$ halves four times ($s = 3$), to $1$;
\item $1 + 1 = 2$, so the next state is $(1,1)$: $3 \cdot 1 + 1 = 4$ halves twice ($s = 1$), back to $1$.
\end{itemize}
So the states run $(1,3)$, $(7,1)$, $(1,2)$, $(1,1)$, and then $(1,1)$ forever.

\paragraph{Nothing is lost (Theorem~\ref{thm:equiv}).} Three facts make the compression faithful. First, $F$ sends the state $(1,1)$ to itself, and the only number with that state is $1$, since $1 + 1 = 2$ gives $m = 1$, $a = 0$ and $\w = 1$. Second, the states an orbit visits, block by block, are exactly the $F$-orbit of its first state (Proposition~\ref{prop:block}(iii)). Third, every \emph{valid} state --- one whose $\w$ is odd, positive and not divisible by $3$, and whose $d$ is at least $1$ --- is the state of some odd number. Put together: an odd number reaches $1$ exactly when its states reach $(1,1)$ (by the first two facts), and every valid state belongs to some odd number (the third). Hence the Collatz conjecture holds if and only if every $F$-orbit from a valid state reaches $(1,1)$. And a cycle of numbers other than $1, 4, 2, 1$ gives a cycle of states other than $(1,1)$ alone, and back: the nontrivial cycles of the two maps correspond one to one. So both of the conjecture's questions can be asked of $F$.

Williams independently found the same states, with their representatives and the forced run of steps, and calls them skeleton rows \cite{williams}; Section~\ref{sec:credit} gives the correspondence.

\section{The exact laws of one step}\label{sec:anchor}

\paragraph{What is claimed.} Once a block is a state $(\w,d)$, three things decide where the orbit goes next: how many halvings, beyond the first, end the block (its exit valuation $s$); how many factors of $3$ the next state picks up (the \emph{absorption}); and the next block's entry depth. The paper gives an exact rule for each (Theorem~\ref{thm:vlaw}, Lemma~\ref{lem:absorption}, Theorem~\ref{thm:depth}). For $s$: sort the states into eight kinds, by the remainder of $\w$ on division by $8$ (which is $1$, $3$, $5$ or $7$, since $\w$ is odd) and by whether $d$ is even or odd. On six kinds $s$ is a fixed number, $1$ or $2$. On the other two --- $\w$ leaving remainder $1$ with $d$ even, or remainder $3$ with $d$ odd, called the \emph{lifting} kinds --- $s$ is at least $3$, and it is $2$ plus the number of lowest binary digits on which $d$ agrees with a fixed binary pattern belonging to $\w$, its \emph{anchor}. A factor of $3$ is gained exactly when $s$ is even. And finitely many binary digits of the state, with three numbers describing the step, either predict the next step exactly or say, correctly, that the next burst of halvings is deep; they never give a wrong answer (Theorems~\ref{thm:deltaM} and~\ref{thm:onestep}).

\paragraph{Why it matters.} One step is completely described. But, by an accounting stated as a heuristic and not proved (Heuristic~\ref{prop:budget}), each decided step uses up binary digits of the starting number, and nothing restores them; this suggests that no finite window of digits can decide an orbit's whole future without further input --- statistical (Section~\ref{sec:aeh}) or about the structure of cycles (Section~\ref{sec:cycles}). In these coordinates the difficulty sits in one place: where the digits of the anchors come from.

\paragraph{A worked example.} For $\w = 1$ the anchor is all zeros, so on its lifting kind ($d$ even) $s$ is $2$ plus the number of trailing zeros of $d$. The example's state $(1,2)$: $d = 2$ is \bin{10} in binary, with one trailing zero, so $s = 3$; and indeed $3 \cdot 5 + 1 = 16$ halves four times, the first halving plus three more. The state $(7,1)$ is of kind ($7$, odd), so $s = 2$; and $s = 2$ is even, so the next state gains a factor of $3$, and indeed $5 + 1 = 6 = 2 \cdot 3$.

\paragraph{The full account.} Appendix~\ref{app:laws} opens with a guided account in plain words --- the table of the eight kinds and why it holds, the factor-of-$3$ rule, the next depth, the window and its three outcomes with their measured frequencies, the digit budget, and the computer checks --- and then gives the statements and proofs.

\section{Cycles: what counting can and cannot do}\label{sec:cycles}

\paragraph{What is claimed.} Suppose a cycle has $p$ blocks; going round it multiplies by $3$ some $n$ times ($n$ is the sum of the blocks' entry depths) and halves some $K$ times. Returning to the start forces $2$ to the power $K$ to be only slightly larger than $3$ to the power $n$, and by less the larger the cycle's numbers are --- because the `add $1$' in each step pushes the number up a little. The closeness is measured by a number $\gamma$: the gap between the two powers is $2^K$ divided by $2^{\gamma}$, so a larger $\gamma$ means a closer approach. Counting (`size') arguments use only this closeness, through the cycle's size conditions (Appendix~\ref{app:cycles}). Classically they give a bound, the uniform \emph{trim} (Rozier 1991; Simons--de~Weger 2005), on how large $n$ can be for $p$ blocks; it weakens exponentially as $p$ grows. Combined with Rhin's theorem --- a classical limit on how close a power of $2$ can come to a power of $3$ --- it leaves only finitely many possible cycles to check at each $p$ (Rozier and Brox proved that there are only finitely many).

The paper does three things here. It rederives that bound by another route. It rederives the classical facts that there are no cycles other than $1, 4, 2, 1$ with one, two or three blocks. And it proves that the exponential weakening cannot be avoided by size arguments. For this it uses an explicit family of \emph{near-cycles} --- lists of block depths and exit valuations shaped like a cycle and required to meet only the size conditions --- the \emph{staircases}. They pass every size test, for every choice of starting block, at every number of blocks from $3$ on (from $16$ on by proof, $3$ to $15$ by computer check; for $2$ and $4$ blocks, passing examples only); and, when every number of the cycle must be at least a given size, at every number of blocks from an explicit threshold on (Theorem~\ref{thm:staircase}).

\paragraph{Why it matters.} A cycle must also meet a second kind of condition: certain whole numbers built from it must divide others. The staircases pass every size test, and every staircase checked fails at least one of these divisibility conditions. So ruling out cycles uniformly in the number of blocks needs divisibility, not size alone --- the conclusion Halbeisen and Hungerb\"uhler reached when measuring cycles by length \cite{hh}. Every statement about the staircases concerns how far size arguments reach; none bears on whether cycles exist.

\paragraph{A worked example.} Take seven blocks with entry depths $4$, $7$, $9$, $15$, $23$, $35$ and $1$, and exit valuations $1$, $1$, $1$, $1$, $1$, $1$ and $49$. Then $n = 94$, the sum of the entry depths, and the halvings total $K = 94 + 55 = 149$ (each block halves once per step, $n$ times in all, plus its exit valuation, $55$ in all). The first six blocks climb; the last falls $49$ extra halvings at once: the profile of an orbit that grows forever, bent into a loop. It passes all seven size tests, yet fails the divisibility conditions, so it is not a cycle. (This instance is v3's, found by an earlier search and stated as Theorem~\ref{thm:staircase}(iv); it is not a member of the family of Theorem~\ref{thm:staircase}(i), whose $\gamma$ lies between about $3.68$ and $5.14$ --- this one's $\gamma$ is $6.74$.)

\paragraph{The full account.} Appendix~\ref{app:cycles} opens with a guided account in plain words --- the closeness measure, with the trivial cycle as an example; the size and divisibility conditions; how known exclusions carry over; the trim's history and the three-step argument that leaves finitely many candidates; this staircase in detail; the sharpness statements and their thresholds; and what is left --- and then gives the statements and proofs.

\section{The fair-coin hypothesis}\label{sec:aeh}

\paragraph{What is claimed.} Read each block as a \emph{letter}, written $m|s$: its entry depth $m$ and its exit valuation $s$, with a bar so that a letter is not mistaken for a state. The classical heuristic picture (Terras \cite{terras}) treats the two parts as independent coin-toss counts --- toss a fair coin until it lands tails, and count the tosses --- so each part is $1$ with chance $\tfrac12$, $2$ with chance $\tfrac14$, and so on. The paper states this as a precise conjecture (Hypothesis~\ref{hyp:aeh}). Draw a large odd starting value at random, and record the letters of its first blocks under two conditions:
\begin{itemize}
\item the \emph{horizon}: the record stops after a number of blocks that is a fixed multiple of the start's number of binary digits;
\item the \emph{budget}: a block's letter is recorded only while the orbit has done fewer halvings than another fixed multiple of that number of digits.
\end{itemize}
Call a choice of the two multiples \emph{admissible} when the record fits inside the budget on the model's own average (four halvings per block) and the budget keeps all but a vanishing share of starts away from the small numbers. The conjecture: for every admissible choice, every finite pattern of letters and every tolerance, the share of starts whose recorded frequency of that pattern differs from its coin frequency by more than the tolerance shrinks to zero as the starting size grows.

Part of it is proved: while the orbit has done fewer halvings than the start has binary digits, the conjecture holds, because the start's lowest binary digits fix its first blocks exactly (the starts with a given pattern of first blocks form one class of remainders on division by a power of $2$). This is the \emph{unconditional base case}, the part proved without assuming anything (proof assembled in the project record, Appendix~\ref{app:verify}). Beyond that it is a conjecture, and its consequences are proved on the condition that it holds.

\paragraph{Why it matters.} Under the coin model the classical facts are exact: for example, a factor of $3$ is gained with chance $\tfrac14 + \tfrac1{16} + \dots = \tfrac13$, because that happens exactly when the exit valuation is even. But they are facts about the model, not about orbits, and no statement that orbits drift downward follows from them. The conjecture says that typical orbits keep following the model after the start's own digits have been used up. It says nothing about where any single orbit ends up.

\paragraph{A worked example.} The running example's blocks $7, 11, 17$; then $13$; then $5$; then $1$ have the letters $3|1$, $1|2$, $1|3$, $1|1$. Under the model the letter $1|2$ has chance $\tfrac12 \times \tfrac14 = \tfrac18$.

\paragraph{The full account.} Appendix~\ref{app:aeh} opens with a guided account in plain words --- the model and what it says by itself; the conjecture, with a worked tally and the meaning of `admissible'; what is proved and why; what the conjecture would and would not give; and the test campaign, whose four apparent anomalies all dissolved under sharper controls --- and then states everything formally.

\section{What remains open}\label{sec:open}

In these coordinates one step is under exact control: the anchor governs it, apart from the factor-of-$3$ absorption, which has its own exact law, and finite windows decide steps without error. What remains is one difficulty seen from two sides: for typical starting values, whether the anchors' digits behave like fair coin tosses (Hypothesis~\ref{hyp:aeh}); for cycles, whether the divisibility conditions of Proposition~\ref{prop:elim} can be met, which size alone cannot decide (Theorem~\ref{thm:staircase}). The staircase belongs to both sides. Three directions are open:
\begin{itemize}
\item \emph{Going beyond size for cycles.} Theorem~\ref{thm:staircase} shows that any uniform exclusion of cycles needs arithmetic statements beyond the size level. One candidate: a cycle's anchors are rigid --- going once round a cycle they return to where they started, so their changes add up to zero.
\item \emph{Following windows step by step.} Study how one exact window of Theorem~\ref{thm:onestep} leads to the next, treating the sequence as a chain (the mathematical method called transfer operators).
\item \emph{The anchors' digits.} Any progress on the digits of the anchors --- in standard terms, digits of $2$-adic logarithms of whole numbers (Appendix~\ref{app:laws}) --- a problem this paper inherits rather than creates.
\end{itemize}

\section{Credit and provenance}\label{sec:credit}

The classical source of each result is named where the result is stated. Related work --- results about almost all starting values (`density results'), comparisons with other work on divisibility by powers of $3$ (`$3$-adic comparisons'), the literature on cycles, and contemporary human--AI work on the problem --- is discussed in Appendix~\ref{app:related}.

The next paragraph says exactly who did what. It uses three terms for objects of Sections~\ref{sec:block}--\ref{sec:reduced}: \emph{skeleton rows}, Williams's name for the states, written $(\lambda,k)$ with $\lambda$ for the odd core $\w$ and $k$ for the depth $d$; the \emph{representative fibre} of a state, the set of numbers that share it; and the \emph{intra-block cascade}, the forced run inside a block (elsewhere in this paper `cascade' means only the block's final burst of halvings). In Williams's picture a row is a pair: a skeleton $\lambda$ (our odd core) and a depth $k$ (our depth).

The division of labor in the present work is stated precisely, because it matters for both credit and blame. The mechanism underneath the reduced coordinates --- each rising step trades a factor of $2$ for a factor of $3$ --- comes from the author's hand-drawn note of March~2023, and is classical (Steiner \cite{steiner}; see also Rozier \cite{rozier90} and LaDue \cite{ladue}). The reduced coordinate itself, the factorization $x + 1 = 2^m 3^a \w$ with its odd core $\w$, was worked out by the author by hand in December~2025, and was named and formalized as the state $(\w,d)$ in AI-assisted work: with ChatGPT (OpenAI) from December~2025, and with Claude (Anthropic) from March~2026. The same states, with their representative fibre and intra-block cascade, appear independently as the skeleton rows $(\lambda,k)$ of Williams \cite{williams}, posted earlier, with $\lambda = \w$ and $k = d$: Williams's Proposition~3.3, Theorem~3.6 and Proposition~3.10 are our Proposition~\ref{prop:block}(iii), (i) and (ii). Williams's open problem~(1) \cite[p.~18]{williams} asks for the map sending one skeleton to the next, and Section~\ref{sec:anchor} answers it in part. When a trajectory leaves a row, Section~\ref{sec:anchor}'s laws give the exit valuation and the entry position, and so the depth $k$ of the row it lands in. They give no law for the skeleton $\lambda$ of that row: on that, $F$ only restates Williams's Proposition~3.10. The analytical superstructure erected on the coordinates --- the anchor formulation, the cycle apparatus, the equidistribution formalization, and the proofs --- was developed in AI-assisted research under the author's direction, drawing on the classical literature's concepts. Appendix~\ref{app:method} states the responsibility and verification protocol.

% =====================================================================
%  APPENDICES
% =====================================================================

\appendix

\section*{Technical summary}
\addcontentsline{toc}{section}{Technical summary}

This is the paper's abstract in the usual notation; the symbols are defined in Appendices~\ref{app:reduced}--\ref{app:aeh} and indexed in Appendix~\ref{app:symbols}.

\begin{quote}\small
The Collatz map sends an odd number $x$ to $3x+1$ with every factor of $2$ divided out. Much of an orbit is bookkeeping: while $x+1$ is divisible by $4$, each step trades one factor of $2$ in $x+1$ for a factor of $3$. We compress each such deterministic valuation run into a single block, producing a self-map $F$ on states $(\w,d)$ --- an odd core prime to $3$ and a depth; the states appear independently as the skeleton rows of Williams (2026) --- whose convergence problem and cycle set are exactly those of the Collatz map. In these coordinates the local arithmetic admits exact laws, governed --- apart from a $3$-adic absorption term, which obeys its own exact law --- by one $2$-adic quantity, the \emph{anchor} $M(\w) = -2\log\w/\log 9 \in \Z_2$: the exit valuation obeys the global law $s = 2 + \vt{d - M(\w)}$ on the lifting classes and is constant on the rest; the depth evolution closes exactly in the anchor displacement together with a stated $3$-adic absorption law; and the anchor increment obeys an exact law modulo any power of $2$, computable from graded residues of the state. A finite window of digits and stratum labels consequently decides each step in an error-free trichotomy, while a digit-budget accounting, stated as a heuristic, indicates that no bounded window can decide infinite horizons --- localizing the difficulty in the digit supply of the anchors.

On the cycle side, a one-line elimination identity (the B\"ohm--Sontacchi cycle equation in block form) yields short rederivations of the classical exclusions at one, two and three blocks and of the classical trim uniform in the number of blocks $p$ (Rozier 1991; Simons--de~Weger 2005), whose constant degrades like $(\LL)^{-p}$. We prove that this exponential loss is intrinsic to size arguments. An explicit family of configurations (\emph{staircases}: divergent-orbit profiles bent into loops) meets every rotation's exact size condition with slack between two absolute constants at every period $p \ge 3$ ($p \ge 16$ by proof, $3 \le p \le 15$ by finite check; at $p \in \{2,4\}$ size-passing configurations are exhibited, with no bracket claimed), and, with every cycle element at least a threshold $X_0$, at every period $p \ge p_1(X_0)$, with slack exceeding the trim's own exponent by at most an absolute constant. At that threshold scale the family realizes the minimax value of Rozier's 1991 extremal under every rotation's exact size condition, the sharpness Simons--de~Weger expected in 2005; the analogue in the cycle-length parameter, and the conclusion that uniform cycle exclusion needs divisibility input, are Halbeisen--Hungerb\"uhler's (1997).

Finally we state the equidistribution (`fair-coin') hypothesis implicit in the classical heuristics as a precise conjecture --- that the block letters of uniformly sampled large starts, tallied within an exponent budget proportional to the start's bit length and for every admissible pair of budget and horizon rates, carry, at every finite block length, the frequencies of an exactly computable Bernoulli law --- prove its consequences conditionally, exhibit an unconditional base case inside the digit budget, and report a calibration campaign whose four apparent anomalies all dissolved under controls.
\end{quote}

\section{Methodology, responsibility, and records}\label{app:method}
This work runs from a hand-drawn origin note (13 March 2023) through seventy-nine drafts to a research wiki whose pages carry status headers, verification records, and refuted claims with their refutation data. It was developed through extensive AI-assisted research in two phases. From December 2025 to March 2026 the author worked with ChatGPT (OpenAI): exploratory definitions, the early CollatzScout drafts, and, alongside them, the author's own hand computations in spreadsheets. From March 2026 the author worked with Claude (Anthropic). The author is responsible for the published claims, the research direction, and the standards of evidence; the AI collaborators carried out the mathematical development within the author's coordinate system --- formulations, proofs, computations, and connections to the literature. Because this division of labor is unusual, correctness claims rest on the documented verification protocol rather than on authorial authority alone: every computational claim is backed by a runnable script, named here or at the record section cited; every result labeled proved was re-verified with independently written code before being so labeled; the calibration campaign's four dissolved anomalies, and two corrected search-filter errors on the cycle side, are documented rather than smoothed over. The complete record (wiki, code, draft history, origin note) is public at \url{https://github.com/macindoe/collatz}; every wiki section and script named in this paper is cited at commit \recordcommit.
% PIN-AT-BUILD: \recordcommit (preamble) must name a commit reachable on the public repository.

\section{The reduced system: definitions and proofs}\label{app:reduced}

This appendix states formally what Sections~\ref{sec:block} and~\ref{sec:reduced} tell in words. Throughout,
\[ T(x) = \frac{3x+1}{2^{\vt{3x+1}}} \qquad (x \text{ odd and positive}) \]
is the odd-to-odd Collatz map (Tao's Syracuse map $\mathrm{Syr}$ \cite{tao}), where $\vt{n}$, the $2$-adic valuation of $n$, is the number of times $2$ divides $n$; and $\vth{n}$ is the number of times $3$ divides $n$. The Collatz conjecture asserts that every $T$-orbit reaches $1$. Each rising step of Section~\ref{sec:block} satisfies $T(x) + 1 = \tfrac32\,(x+1)$ (proof of Proposition~\ref{prop:block}(i)).

\begin{definition}\label{def:reduced}
A \emph{reduced state} is a pair $(\w, d)$ with $\w$ odd and positive, $3\nmid\w$, $d \ge 1$. Set
\[ A = 3^d\w - 1,\qquad s = \vt{A},\qquad x_{\mathrm{exit}} = A/2^s,\qquad C = A + 2^s, \]
\[ \sigma = \vt{C},\qquad a_+ = \vth{C},\qquad m_+ = \sigma - s,\qquad \wnext = \frac{C}{2^{\sigma}3^{a_+}},\qquad \dnext = m_+ + a_+, \]
and define the \emph{reduced map} $F(\w,d) = (\wnext, \dnext)$. The \emph{projection} $R$ sends an odd $x$ with $x + 1 = 2^m 3^a \w$ ($3\nmid\w$) to $R(x) = (\w, m + a)$.
\end{definition}

\begin{proposition}[block structure and well-definedness]\label{prop:block}
Let $x$ be odd with $x + 1 = 2^m u$, $u = 3^a\w$ odd, $3 \nmid \w$, and let $(\w,d) = R(x)$, $d = m+a$. Then:
\begin{enumerate}
\item[(i)] the next $m - 1$ iterates of $T$ are $x_j$ with $x_j + 1 = 2^{m-j}\,3^j u$ for $j = 0,\dots,m-1$;
\item[(ii)] the following iterate is $T(x_{m-1}) = A/2^s = x_{\mathrm{exit}}(\w,d)$, which depends only on $(\w,d)$;
\item[(iii)] consequently $R(T^{m}(x)) = F(R(x))$ for every representative $x$ of $(\w,d)$, the representatives being exactly $x = 2^{d-a}3^a\w - 1$, $0 \le a \le d-1$; and the sequence of reduced states visited by any $T$-orbit is the $F$-orbit of its projection.
\end{enumerate}
\end{proposition}

\begin{proof}
(i) If $y + 1 = 2^k v$ with $k \ge 2$, $v$ odd, then $3y + 1 = 3(y+1) - 2 = 2\left(3\cdot 2^{k-1}v - 1\right)$, an odd multiple of $2$, so $T(y) = 3\cdot2^{k-1}v - 1$ and $T(y) + 1 = 2^{k-1}(3v)$: one factor of $2$ is exchanged for one factor of $3$. Induction from $y = x$ gives (i). (ii) At $j = m-1$ we have $x_{m-1} + 1 = 2\cdot 3^{m-1}u$, so $3x_{m-1} + 1 = 3(x_{m-1}+1) - 2 = 2(3^m u - 1) = 2A$ since $3^m u = 3^{m+a}\w = 3^d \w$. Thus $T(x_{m-1}) = 2A/2^{1+s} = A/2^s$, and $A$ depends only on $(\w,d)$. (iii) $x_{\mathrm{exit}} + 1 = (A + 2^s)/2^s = C/2^s$, so $R(x_{\mathrm{exit}}) = (\wnext, (\sigma - s) + a_+) = F(\w,d)$ by Definition~\ref{def:reduced}; the representative list is immediate from $x + 1 = 2^{d-a}3^a\w$.
\end{proof}

\begin{theorem}[convergence translation]\label{thm:equiv}
$(1,1)$ is a fixed point of $F$ whose $R$-fiber is $\{1\}$. The Collatz conjecture holds if and only if every $F$-orbit from a valid state reaches $(1,1)$, and $R$ induces a bijection between nontrivial $T$-cycles and $F$-cycles other than $\{(1,1)\}$.
\end{theorem}

\begin{proof}
For $(1,1)$: $A = 2$, $s = 1$, $x_{\mathrm{exit}} = 1$, $C = 4$, so $F(1,1) = (1,1)$; the fiber is $\{1\}$ since the representatives are $x = 2\cdot 1 - 1 = 1$ only. If the conjecture holds and $(\w,d)$ is valid, its representative $x = 2^d\w - 1$ reaches $1$, so by Proposition~\ref{prop:block}(iii) the $F$-orbit of $(\w,d)$ reaches $R(1) = (1,1)$. Conversely if every $F$-orbit reaches $(1,1)$, then any odd $x$ has a $T$-orbit visiting the fiber of $(1,1)$, i.e.\ the value $1$. For cycles: a nontrivial $T$-cycle projects to a periodic $F$-orbit avoiding $(1,1)$ (its values avoid $1$); conversely an $F$-cycle's block entries $e_t = x_{\mathrm{exit}}(S_{t-1})$ are deterministic in the states and joined by genuine $T$-iteration, so $e_0$ is $T$-periodic and nontrivial; the two constructions are mutually inverse.
\end{proof}

\section{The exact laws: statements and proofs}\label{app:laws}

This appendix states formally what Section~\ref{sec:anchor} tells in words. It opens with a guided account in plain words, using the running example of Section~\ref{sec:reduced}; the formal statements follow.

\subsection*{A guided account}

Once a block is a state $(\w,d)$, three questions decide where the orbit goes next: how many halvings, beyond the first, end the block (its exit valuation $s$); how many factors of $3$ the next state picks up (the \emph{absorption}); and the next block's entry depth. Each has an exact answer, and together they give the next state. The running example's states $(1,3)$, $(7,1)$, $(1,2)$, $(1,1)$ illustrate each law.

\paragraph{The valuation law (Theorem~\ref{thm:vlaw}).} The rule for $s$ depends first on which of eight kinds a state is. Sort the states by the remainder of $\w$ on division by $8$ (which is $1$, $3$, $5$ or $7$) and by whether $d$ is even or odd:
\begin{center}
\begin{tabular}{c|cccc}
\hline
 & \multicolumn{4}{c}{remainder of $\w$ on division by $8$} \\
 & $1$ & $3$ & $5$ & $7$ \\
\hline
$d$ even & lifting & $s = 1$ & $s = 2$ & $s = 1$ \\
$d$ odd & $s = 1$ & lifting & $s = 1$ & $s = 2$ \\
\hline
\end{tabular}
\end{center}
On the six kinds marked with a number, $s$ is that number. On the two \emph{lifting} kinds, $s$ is at least $3$, and it is set by a fixed binary pattern that belongs to $\w$, running on to the left without end: $s$ is $2$ plus the number of lowest binary digits on which $d$ agrees with that pattern, counted from the right up to the first disagreement.

For $\w = 1$ the pattern is all zeros, so on its lifting kind ($d$ even) $s$ is $2$ plus the number of trailing zeros of $d$:
\begin{center}
\begin{tabular}{rrrr}
\hline
$d$ & $d$ in binary & trailing zeros & $s$ \\
\hline
$2$ & \bin{10} & $1$ & $3$ \\
$4$ & \bin{100} & $2$ & $4$ \\
$6$ & \bin{110} & $1$ & $3$ \\
$8$ & \bin{1000} & $3$ & $5$ \\
\hline
\end{tabular}
\end{center}
The state $(1,2)$ of the example is the first row: $3 \cdot 5 + 1 = 16$ halves four times, the first halving plus $s = 3$ more. The example's other states are not lifting: $(1,3)$ and $(1,1)$ are of kind ($1$, odd), so $s = 1$; $(7,1)$ is of kind ($7$, odd), so $s = 2$. These match the halvings counted in Section~\ref{sec:reduced}.

Why the table holds. The cascade begins at the number $2 \times (3^d\w - 1)$ (Appendix~\ref{app:reduced}), so $s$ is the number of times $2$ divides $3^d\w - 1$; for $(1,3)$, $27 - 1 = 26 = 2 \cdot 13$, so $s = 1$. The powers of $3$ leave remainders $3, 1, 3, 1, \dots$ on division by $8$, so the kind fixes the remainder of $3^d\w$. On the lifting kinds that remainder is $1$, so $3^d\w - 1$ is divisible by $8$ and $s$ is at least $3$. On the other six it is $3$, $5$ or $7$, so $3^d\w - 1$ leaves remainder $2$, $4$ or $6$, and $s$ is $1$ or $2$.

The pattern is called the \emph{anchor} of $\w$; Definition~\ref{def:anchor} defines it. The law describes a step; it is not offered as a faster way to compute one.

\paragraph{The factor of $3$ (Lemma~\ref{lem:absorption}).} \emph{Whether} the next state gains a factor of $3$ is simple. One more than the exit, multiplied by $2^s$, is the number $3^d\w + 2^s - 1$, so the next state's factors of $3$ are those of this number. Its first part is a multiple of $3$, and $2^s - 1$ is a multiple of $3$ exactly when $s$ is even (the powers of $2$ leave remainders $2, 1, 2, 1, \dots$ on division by $3$). So the next state gains a factor of $3$ exactly when $s$ is even. In the example, the block of $13$, state $(7,1)$, has $s = 2$; its exit is $5$, and $(5 + 1) \times 4 = 24 = 21 + 3$, a multiple of $3$, so the next state $(1,2)$ has $a = 1$. The blocks with $s = 1$ and $s = 3$ gain none. \emph{How many} factors of $3$ are gained is given by an exact formula; in one boundary case it depends on the remainders of $\w$ on division by powers of $3$.

\paragraph{The next depth (Theorem~\ref{thm:depth}).} The next state's depth is the next block's entry depth plus the absorption. The absorption is given by Lemma~\ref{lem:absorption}. The next entry depth is given by an exact formula in $s$ and in how $d$ compares with the anchor (Theorem~\ref{thm:depth}); the formula can never give less than $1$. In the example, from $(7,1)$ the next number is $5$, and $5 + 1 = 6$ has one factor of $2$ and one of $3$: entry depth $1$, absorption $1$, next depth $2$.

\paragraph{How many digits one step needs (Theorem~\ref{thm:deltaM}).} The anchor never ends, but one step needs only finitely many of its digits: $s$ depends only on how many of its lowest digits agree with $d$. The paper makes this exact for the anchor itself. Fix how many of the lowest binary digits of the change in anchor, from one state to the next, are wanted. Then finitely many binary digits of $\w$ and of $d$, together with the lowest digits of the absorption, determine them, by an explicit computation. The same digits of $\w$ and $d$ also determine $s$ and the next entry depth. They do not determine the absorption, which depends on the remainders of $\w$ on division by powers of $3$; so the absorption has to be supplied separately, as a label of the step.

\paragraph{Three outcomes, never wrong (Theorem~\ref{thm:onestep}).} Call that finite data --- the digits of $\w$ and $d$, together with the step's three labels ($s$, the next entry depth and the absorption) --- a \emph{window}. The \emph{depth} of a window is the number of anchor digits it tracks. From the window alone:
\begin{itemize}
\item the next depth and the next state's kind are determined exactly;
\item if the next kind is not lifting, the next $s$ ($1$ or $2$) is determined exactly;
\item if it is lifting, either the next $s$ is determined exactly, or the window reports, correctly, that the next $s$ is \emph{deep}: at least $2$ more than the window's depth.
\end{itemize}
The window never gives a wrong value; it is undecided only in the deep case. How often is that? Under the fair-coin model of Section~\ref{sec:aeh} (whose cut-off, the cap, is taken at least the window's depth plus $2$), exactly $1$ time in $2$ raised to the window's depth plus $1$: $1$ in $32$ ($3.125\%$) at window depth $4$, and $1$ in $512$ ($0.195\%$) at depth $8$. Along real orbits the record measures $2.75\%$ at depth $4$ and $0.19\%$ at depth $8$.

\paragraph{The digit budget (Heuristic~\ref{prop:budget}; a heuristic, not a theorem).} Each decided step uses binary digits of the state to a depth a little beyond the window's own: the window's depth, plus $s$, plus the next entry depth, plus $2$ (under the model of Section~\ref{sec:aeh}, $s$ plus the next entry depth averages $4$). Theorem~\ref{thm:deltaM} shows that these digits suffice; that they are also needed is part of the heuristic. Nothing restores the digits used: the anchors of later states are determined by the full initial state, not by any finite window of it. So, by this count --- about the window's depth plus $6$ digits per step on average --- a window of finitely many digits can decide at most a number of steps proportional to its length. This strongly suggests --- and the paper treats it as its organizing heuristic, not as a proved theorem --- that a finite window cannot decide an orbit's whole future without additional input: statistical input (Section~\ref{sec:aeh}), or input about the rigid structure of cycles (Section~\ref{sec:cycles}).

\paragraph{Checks.} The valuation, entry-depth and absorption laws were checked by independently written code, with zero discrepancies: the valuation law on $8{,}000$ random states, the entry-depth law on all $63{,}000$ valid states with $\w$ below $3000$ and $d$ below $64$, $\w = 1$ included (\texttt{experiments/depth\_law\_reverify.py}), and the absorption law on $327{,}980$ states (Remark~\ref{rem:verify1}).
The three outcomes were run along real orbits, $21{,}296$ steps, with zero errors (Remark~\ref{rem:verify2}). Appendix~\ref{app:verify} lists the scripts.

\paragraph{Where the difficulty is.} Every quantity above except the absorption obeys an exact law in how $d$ compares with the anchor, and the absorption obeys its own exact law (Lemma~\ref{lem:absorption}). A finite window decides the next step in three outcomes and is never wrong. And the digit-budget accounting, stated as a heuristic, indicates that each decision uses up digits that are not restored. So in these coordinates the difficulty of the problem sits in one place: it is the question of where the digits of the anchors come from. Heuristic~\ref{prop:budget} is the paper's organizing negative observation. It splits what remains into a statistical question for typical starting values (Section~\ref{sec:aeh}) and a question about the rigid arithmetic of cycles (Section~\ref{sec:cycles}); we know of nothing remaining outside those two.

\subsection*{The formal statements}

For $u \equiv 1 \pmod 8$ the equation $9^{t} = u^{-1}$ has a unique solution $t \in \Z_2$, since $9$ topologically generates $1 + 8\Z_2$; write $N(u)$ for it, so $N(u) = -\log u/\log 9$ in the $2$-adic logarithm. $N$ is an isomorphism $(1+8\Z_2,\times)\to(\Z_2,+)$ --- the inverse of $t \mapsto 9^{-t}$, a continuous homomorphism, injective by \eqref{eq:isometry} below and onto because its image is compact and dense --- and
\begin{equation}\label{eq:isometry}
\vt{9^t - 1} = 3 + \vt{t}\qquad (t \in \Z_2\setminus\{0\}),
\end{equation}
the \emph{isometry}, since $9 = 1 + 8$ and $\vt{\log(1+8z)} = 3 + \vt{z}$-type estimates hold with equality here.

\begin{definition}[anchor]\label{def:anchor}
For every odd $\w$ set $M(\w) = N(\w^2) = -2\log\w/\log 9 \in \Z_2$, well defined since $\w^2\equiv 1 \pmod 8$. $M(\w)$ is even iff $\w\equiv\pm1\pmod 8$.
\end{definition}

The parity clause: $\w^2 \equiv 1 \pmod{16}$ for $\w \equiv \pm 1 \pmod 8$ and $\w^2 \equiv 9 \pmod{16}$ otherwise, while $9^t \equiv \w^{-2} \pmod{16}$ and $9^2 \equiv 1 \pmod{16}$ force $t = M(\w)$ even in the first case and odd in the second.

\begin{lemma}[squaring lemma]\label{lem:squaring}
Let $\w$ be odd, $d \ge 1$, and $Y = 3^d\w$. If $Y \equiv 1 \pmod 8$, then
\[ \vt{Y - 1} = 2 + \vt{d - M(\w)}. \]
\end{lemma}

\begin{proof}
$Y^2 = 9^d\w^2 = 9^{d}\,9^{-M(\w)}\cdot\bigl(9^{M(\w)}\w^2\bigr) = 9^{\,d - M(\w)}$, since $9^{M(\w)} = \w^{-2}$ by definition of $N$. By \eqref{eq:isometry}, $\vt{Y^2 - 1} = 3 + \vt{d - M(\w)}$ (nonzero since $Y^2 \neq 1$ for a valid state). Since $Y \equiv 1 \pmod 8$ we have $\vt{Y+1} = 1$, and $\vt{Y-1} = \vt{Y^2-1} - \vt{Y+1}$ gives the claim.
\end{proof}

\begin{theorem}[global valuation law]\label{thm:vlaw}
Call $(\w,d)$ \emph{lifting} if $3^d\w \equiv 1 \pmod 8$, i.e.\ $\w\equiv1\ (8)$ with $d$ even or $\w\equiv3\ (8)$ with $d$ odd. On the lifting classes
\[ s = \vt{3^d\w - 1} = 2 + \vt{d - M(\w)} \;\ge\; 3, \]
while on the six non-lifting residue-parity classes $s \in \{1,2\}$ is the constant determined by $(\w \bmod 8,\, d \bmod 2)$: $s = 1$ on $(1,\mathrm{odd}), (5,\mathrm{odd}), (7,\mathrm{even}), (3,\mathrm{even})$ and $s = 2$ on $(5,\mathrm{even}), (7,\mathrm{odd})$. Unconditionally $s \le C(\w)\,(1 + \log d)^2$: for $\w \ne 1$ by $p$-adic Baker theory \cite{yu}, \cite[Corollaire~2]{bl}, and at $\w = 1$ exactly, since there the lifting class is $d$ even and $s = 2 + \vt{d}$.
\end{theorem}

\begin{proof}
The lifting case is Lemma~\ref{lem:squaring}; the parity claim ($d \equiv M(\w) \bmod 2$ on lifting classes) follows from $\vt{d - M(\w)} \ge 1$ there. Off-lifting, $3^d\w \bmod 8 \in \{3,5,7\}$ gives $\vt{3^d\w - 1} \in \{1,2\}$ by inspection of residues mod $8$. The Baker bound applies to $s = \vt{9^n\tilde\w - 1}$, a linear form in two $2$-adic logarithms, for $\tilde\w \ne 1$ (here $\tilde\w \in \{\w, 3\w\}$ according to the lifting component, so $\tilde\w = 1$ only at $\w = 1$, $d$ even); at $\tilde\w = 1$ the pair $(9,1)$ is multiplicatively dependent and the Baker input does not apply, but $M(1) = 0$ and the lifting case above gives $s = 2 + \vt{d} \le 2 + \log_2 d$ exactly, by the isometry \eqref{eq:isometry}.
\end{proof}

\begin{lemma}[$3$-adic absorption law]\label{lem:absorption}
Let $h(s) = \vth{2^s - 1}$, so $h(s) = 0$ for $s$ odd and $h(s) = 1 + \vth{s}$ for $s$ even. Then
\[ a_+ = \vth{C} = \begin{cases} 0, & s \text{ odd},\\[2pt] \min(d,\, h(s)), & s \text{ even},\ d \ne h(s),\\[2pt] d + \vth{\w + (2^s-1)3^{-d}}, & s \text{ even},\ d = h(s). \end{cases} \]
In particular the next state gains a factor of $3$ exactly when $s$ is even.
\end{lemma}

\begin{proof}
$C = 3^d\w + (2^s - 1)$ with $\vth{3^d\w} = d$ exactly and $\vth{2^s-1} = h(s)$; $h(s) = 0$ iff $s$ odd since $2^s \equiv 2 \pmod 3$ then, and for even $s$, lifting-the-exponent gives $\vth{4^{s/2}-1} = \vth{3} + \vth{s/2} = 1 + \vth{s}$. The ultrametric equality $\vth{X+Y} = \min(\vth X,\vth Y)$ holds when the valuations differ; when $d = h(s)$, factor $3^d$ and evaluate the remaining unit sum.
\end{proof}

\begin{theorem}[depth closure]\label{thm:depth}
On every residue class the next depth is $\dnext = m_+ + a_+$ with $a_+$ given exactly by Lemma~\ref{lem:absorption} and $m_+$ by the \emph{entry-depth law}: writing $\xi$ for the anchor displacement ($\xi = n - N(\w)$ on the lifting class $d = 2n$, $\xi = n - N(3\w)$ on $d = 2n+1$) and $\delta_s = N(1 - 2^s)$ (which satisfies $\vt{\delta_s} = s - 3$ for $s \ge 3$),
\[ m_+ = 3 - s + \vt{\xi + \delta_s} \quad\text{(lifting)},\qquad
   m_+ = \begin{cases} 1 + \vt{d - M(\w)}, & s = 1,\\ \vt{(d-1) - M(\w)}, & s = 2. \end{cases} \]
The forced carry $\vt{\delta_s} = s-3 = \vt{\xi}$ yields $m_+ \ge 1$ structurally.
\end{theorem}

\begin{proof}
All cases reduce to the \emph{target-shift lemma}: for $c \in 1 + 8\Z_2$ and $d = 2n$ with $\w \equiv 1\ (8)$,
\[ \vt{3^{d}\w - c} \;=\; 3 + \vt{n - N(\w) + N(c)}, \]
proved exactly as Lemma~\ref{lem:squaring} after writing $3^d\w - c = 9^{-N(c)}\bigl(9^{\,n - N(\w) + N(c)} - 1\bigr)$. Since $m_+ = \vt{x_{\mathrm{exit}}+1}$ and $2^s(x_{\mathrm{exit}}+1) = C = 3^d\w - (1 - 2^s)$, the lifting case is the shift $c = 1 - 2^s$ (which lies in $1+8\Z_2$ for $s \ge 3$), with $\vt{\delta_s} = s - 3$ from the isometry; the odd lifting component reduces to the even one via the companion parameter $3\w$. For $s = 1$: the shifted target is $-1$, and $\vt{X+1} = \vt{X^2-1} - \vt{X-1}$ with $X = 3^d\w$, $X^2 = 9^{\,d - M(\w)}$, gives the formula; For $s = 2$: here $C = 3^d\w + 3 = 3(X' + 1)$ with $X' = 3^{d-1}\w$, so $m_+ = \vt{C} - 2 = \vt{X'+1} - 2$. The $s=2$ classes force $3^d\w \equiv 5 \pmod 8$, hence $X' \equiv 3^{-1}\cdot 5 \equiv 7 \pmod 8$, so $\vt{X'-1} = 1$; and $X'^2 = 9^{\,(d-1) - M(\w)}$ as in Lemma~\ref{lem:squaring}, so $\vt{X'+1} = \vt{X'^2-1} - \vt{X'-1} = \bigl(3 + \vt{(d-1) - M(\w)}\bigr) - 1$, giving $m_+ = \vt{(d-1) - M(\w)}$ as claimed.
\end{proof}

\begin{remark}[verification notes]\label{rem:verify1}
Each law above was checked by independently written code, with zero discrepancies: the valuation law on $8{,}000$ random states across both lifting components, the entry-depth law exhaustively on all $63{,}000$ valid states with $\w < 3000$, $1 \le d < 64$ ($\w = 1$ included, where every law holds as stated), and the absorption law on $327{,}980$ states ($983{,}954$ comparisons against both this Lemma and its form in the record), with the boundary branch $d = h(s)$, on which $a_+$ is a $3$-adic function of $\w$, and the shallow branch $d < h(s)$ constructed rather than sampled. The records are \texttt{stage3.md} \S11.8.6.3 and that page's status header; the code is \texttt{experiments/depth\_law\_reverify.py} and \texttt{experiments/absorption\_law.py}.
\end{remark}

\begin{theorem}[low-order anchor-increment law]\label{thm:deltaM}
Let $\Delta M = M(\wnext) - M(\w)$. For every $k \ge 1$, $\Delta M \bmod 2^k$ is determined by, and explicitly computable from,
\[ \w \bmod 2^{\sigma + k + 2},\qquad d \bmod 2^{\sigma + k},\qquad a_+ \bmod 2^k, \]
and the stratum data $(s, \sigma)$ are determined by the same residues. This last clause does not extend to $a_+$: by Lemma~\ref{lem:absorption} the absorption is a $3$-adic function of $\w$, and it enters the display only through its residue $a_+ \bmod 2^k$.
\end{theorem}

\begin{proof}
Three digit-determinacy facts chain together. (a) For $u \equiv 1\ (8)$, $N(u) \bmod 2^k$ depends only on $u \bmod 2^{k+3}$, since $9$ has order $2^k$ mod $2^{k+3}$ and $9^t \equiv u^{-1}$ determines $t$ mod $2^k$. (b) $C \bmod 2^{q}$ is determined by $\w \bmod 2^q$, $d \bmod 2^{q-2}$, $s$, since $3$ has order $2^{q-2}$ mod $2^q$. (c) $\wnext \bmod 2^r$ is determined by $C \bmod 2^{\sigma + r}$, $\sigma$, and $a_+ \bmod 2^{r-2}$, by exact division and the order of $3$ mod $2^r$. Now $\Delta M = N\bigl((\wnext/\w)^2\bigr)$; by (a) this needs $(\wnext/\w)^2 \bmod 2^{k+3}$, hence $\wnext, \w \bmod 2^{k+2}$; by (c) with $r = k+2$ this needs $C \bmod 2^{\sigma+k+2}$ and $a_+ \bmod 2^k$; by (b) this needs the stated residues of $\w$ and $d$. Finally $s \le \sigma$ and $\sigma$ are visible in $A \bmod 2^{\sigma+1}$, $C \bmod 2^{\sigma+1}$.
\end{proof}

\begin{theorem}[one-step trichotomy]\label{thm:onestep}
The \emph{depth-$k$ window} of a valid state $(\w,d)$ consists of the residues of Theorem~\ref{thm:deltaM} together with the stratum labels $(s, \sigma, a_+)$. From the depth-$k$ window alone, one can determine: the exact next depth $\dnext$ and next residue-parity class; the exact value $s_+ \in \{1,2\}$ if the next class is non-lifting; and, if lifting, the truncation of the next displacement mod $2^k$ --- yielding $s_+$ exactly when nonzero, and the correct report ``$s_+ \ge k+2$'' when zero. The window never outputs a wrong value.
\end{theorem}

\begin{proof}
By facts (b), (c) above the window determines $\wnext \bmod 2^{k+2}$, hence the next class; $\dnext = (\sigma - s) + a_+$ is exact from the stratum data. Non-lifting $s_+$ is the class constant (Theorem~\ref{thm:vlaw}). Lifting: $M(\wnext) \bmod 2^k$ follows from fact (a) applied to $\wnext^2$, so $(\dnext - M(\wnext)) \bmod 2^k$ is computable; if nonzero its valuation is visible and Theorem~\ref{thm:vlaw} gives $s_+$; if zero then $\vt{\dnext - M(\wnext)} \ge k$, i.e.\ $s_+ \ge k + 2$.
\end{proof}

By the proof, the window is undecided exactly when $s_+ \ge k+2$. Under the law $\pi_{k,D}$ of Appendix~\ref{app:aeh} with $D \ge k+2$, whose $s$-marginal is $P(s = j) = 2^{-j}$ for $j < D$, that event has probability $\sum_{j \ge k+2} 2^{-j} = 2^{-(k+1)}$ exactly. Along real orbits the record measures an undecided rate of $0.0275$ against $0.03125$ at $k = 4$, and $0.0019$ against $0.00195$ at $k = 8$ (\texttt{stage4.md}, Remark ``undecided rate'').

\begin{heuristic}[digit budget]\label{prop:budget}
Deciding one step at depth $k$ consumes the state's $2$-adic data to depth $\sigma + k + 2$ ($\sigma = s + m_+$ has mean $4$ under the model law of Appendix~\ref{app:aeh}, where $s$ and $m_+$ are geometric of mean $2$; the heuristic does not depend on the constant). Theorem~\ref{thm:deltaM} shows that these residues suffice; that they are also needed is part of the heuristic. Nothing regenerates them: the anchors of later states are functions of the full initial state, not of any bounded window. A window of $W$ digits therefore supports $O(W/(k+6))$ decided steps. This strongly suggests --- and we treat it as the organizing heuristic, not a formalized theorem --- that bounded-window determinism cannot decide unbounded horizons without additional statistical or rigidity input.
\end{heuristic}

\begin{remark}[verification note]\label{rem:verify2}
The trichotomy was run along real orbits for $k \in \{4,8\}$: $21{,}296$ steps, $21{,}000$ decided with zero errors, $296$ flagged deep with zero violations of the bound $s_+ \ge k+2$. Code: \texttt{experiments/one\_step\_propagation.py}; the low-order law of Theorem~\ref{thm:deltaM} is checked by \texttt{experiments/anchor\_increment.py}.
\end{remark}

\section{Cycles: the apparatus and the proofs}\label{app:cycles}

This appendix states formally what Section~\ref{sec:cycles} tells in words. It opens with a guided account in plain words; the formal statements follow.

\subsection*{A guided account}

A counterexample to the conjecture would be either an orbit that grows forever or an orbit that falls into a cycle not containing $1$. Cycles are the one case where finitely many digits are enough: a cycle is a finite list of numbers, so arithmetic on that list can settle it, without the unending supply of digits that Section~\ref{sec:anchor} says a general orbit needs. This account explains what the classical counting (`size') arguments achieve against cycles, and why they go no further: for every number of blocks from $3$ on, and, for every size threshold, from an explicit threshold number of blocks on, the size conditions are met at the scale the classical bound allows. (That bound is called the \emph{trim}; it is explained below, in ``The uniform trim''.) Two examples run through it: the trivial cycle $1, 4, 2, 1$ and a staircase of seven blocks.

\paragraph{A cycle as numbers.} Suppose $F$ has a cycle of $p$ blocks. Write $n$ for the total of the blocks' entry depths, and $K$ for the total number of halvings around the cycle. A block of entry depth $m$ multiplies by $3$ exactly $m$ times (once per step), and halves $m$ plus $s$ times (once in each of its $m - 1$ rising steps, then $1 + s$ times in its cascade). So $n$ counts the multiplications by $3$, and $K$ is $n$ plus the total of the exit valuations. In the trivial cycle, the state $(1,1)$ is a single block ($p = 1$) with $m = 1$ and $s = 1$: multiply by $3$ once and halve twice, so $n = 1$ and $K = 2$.

\paragraph{Known exclusions carry over.} In the one-division map of Section~\ref{sec:block}, a cycle of $p$ blocks has exactly $p$ runs of consecutive odd numbers, and its $p$ local minima (numbers smaller than both their neighbours) are exactly its block entries. So a cycle with $p$ blocks is a $p$-cycle in Simons and de~Weger's sense (their `$m$-cycle', with their $m$ equal to our $p$, not to an entry depth) \cite{sdw}, and $p$ is what Rozier calls the number of local minima \cite{rozier91}; results stated in their terms apply here directly. For example, Hercher's theorem \cite{hercher} rules out every nontrivial cycle of $F$ (every cycle other than $1, 4, 2, 1$) with at most $91$ blocks. Such results bound the number of blocks only; they say nothing about how many blocks have entry depth $2$ or more.

\paragraph{Size.} Going round the cycle multiplies by $3$ exactly $n$ times and halves exactly $K$ times, and the `add $1$' in every step pushes the number up a little each time. Multiplying by $3^n$ and dividing by $2^K$ alone would have to land back on the start; with the extra $1$s pushing up, the halvings must win slightly. Returning to the start therefore forces $2$ to the power $K$ to exceed $3$ to the power $n$ --- but only by a little, and by less the larger the cycle's numbers are (Corollary~\ref{cor:size}). Write $q$ for the excess. In the trivial cycle, $2^2 = 4$ exceeds $3^1 = 3$ by $q = 1$. \emph{Size arguments} are arguments from how close $2^K$ must be to $3^n$. A convenient measure of that closeness is $\gamma$, defined by: $3^n$ equals $2^K$ times $(1 - 1/2^{\gamma})$. The larger $\gamma$, the closer the two powers. In the trivial cycle, $3 = 4 \times (1 - 1/4)$, so $\gamma = 2$.

\paragraph{Size and divisibility (Proposition~\ref{prop:elim}).} The \emph{elimination identity} (the cycle equation of B\"ohm and Sontacchi \cite{bohm}, written in blocks) turns a cycle into $p$ conditions, one for each choice of starting block, or \emph{rotation}. Each says: an explicit positive whole number, built from the cycle's entry depths and exit valuations, equals $q$ times a positive whole number. Two kinds of condition follow. A positive whole multiple of $q$ is at least $q$, so $q$ is at most each of these $p$ numbers: the \emph{size conditions}. And $q$ divides each of them: the \emph{divisibility conditions}. In the trivial cycle the one number is $4 - 3 = 1 = q$, and both conditions hold. Any list of positive entry depths and exit valuations for which $2^K$ exceeds $3^n$ is a \emph{configuration}; one meeting every size condition is a \emph{size-passer}. A cycle gives a size-passer, but a size-passer need not give a cycle: it may fail the divisibility conditions, as the staircases below do.

\paragraph{Small periods (Theorem~\ref{thm:smallp}).} $F$ has no nontrivial cycles of $1$, $2$ or $3$ blocks. These statements are contained in the classical results \cite{steiner,rozier90,sdw,hercher}; the derivations, not the theorems, are the contribution.

\paragraph{The uniform trim: a short history.}
\begin{itemize}
\item In 1991 Rozier bounded, for each number of local minima of a cycle (our $p$), the number of its odd steps (our $n$), by a bound whose strength fades geometrically as $p$ grows; with it he proved that only finitely many cycles have a given $p$ (his Th\'eor\`eme). He observed that the bound is the best possible if the block depths may be any real numbers rather than whole numbers; the best real depths grow geometrically, by a factor of $1.585\ldots$ (the base-$2$ logarithm of $3$) from each block to the next \cite{rozier91}.
\item In 1997 Halbeisen and Hungerb\"uhler, measuring cycles by their length rather than their number of blocks, proved the size criterion optimal: cycles of rational numbers of every length attain it. They reduced the remaining question to divisibility \cite{hh}.
\item In 2000 Brox proved again that only finitely many cycles can have any given number of local minima \cite{brox}.
\item In 2005 Simons and de~Weger made Rozier's bound explicit, computed with it, and expected its exponential dependence on the number of minima to be unavoidable \cite{sdw}.
\end{itemize}
In this paper's notation, and up to an additive constant, the bound --- the \emph{uniform trim} --- says (Theorem~\ref{thm:uniform}): in every nontrivial cycle with $p$ blocks, $\gamma$ plus the base-$2$ logarithm of $p$ is larger than $n$ times a factor that shrinks like $1.585\ldots$ to the power $-p$. Read it this way: the more multiplications by $3$ a cycle has (the larger $n$), the closer $2^K$ must be to $3^n$; but because the factor shrinks, that requirement weakens exponentially as the number of blocks grows. The proof uses only the size conditions, so the bound holds for every size-passer, not only for cycles.

How this leaves finitely many candidates at each number of blocks, in three steps (Theorem~\ref{thm:uniform}, second part):
\begin{enumerate}
\item[(a)] By the trim, a cycle with $p$ blocks can have a large $n$ only if its $\gamma$ is large --- only if $2^K$ is extremely close to $3^n$.
\item[(b)] By Rhin's theorem \cite{rhin}, a power of $2$ cannot come too close to a power of $3$: $\gamma$ can grow at most like a fixed multiple of the logarithm of $n$.
\item[(c)] A quantity growing in proportion to $n$ cannot stay below one growing like the logarithm of $n$ for ever. So, for each $p$, $n$ is bounded --- at most of the order of $p$ times ($1.585\ldots$ raised to the power $p$) --- and only finitely many configurations remain to check.
\end{enumerate}
Our proof is a rederivation by another route, chosen because it identifies the profile, extremal for the bound, that the staircase realizes.

\paragraph{The staircase, as a picture (Theorem~\ref{thm:staircase}(iv)).} Take seven blocks with entry depths $4$, $7$, $9$, $15$, $23$, $35$ and $1$, and exit valuations $1$, $1$, $1$, $1$, $1$, $1$ and $49$. Then $p = 7$; the entry depths total $n = 94$; the exit valuations total $55$; and so the halvings total $K = 94 + 55 = 149$. The first six blocks climb, their depths growing roughly geometrically, each ending with the smallest possible burst ($s = 1$); the last block has depth $1$ and $s = 49$, one long fall. This is the profile of an orbit that grows forever, bent into a loop by one crash. This seven-block instance is v3's, found by an earlier search and stated as part~(iv) of Theorem~\ref{thm:staircase}; it is not one of the family of part~(i) below, whose $\gamma$ lies between about $3.68$ and $5.14$.

It passes all seven size tests, with $\gamma = 6.74$. The bound known for three-block cycles (Theorem~\ref{thm:staircase}(iv) states it) demands that $\gamma$ exceed $0.1157$ times $n$, minus $2$: about $8.9$ at $n = 94$. So a bound as strong as the three-block one would wrongly exclude this seven-block size-passer: at seven blocks the requirement must be weaker than at three. (That it must weaken exponentially as the number of blocks grows is part~(i) below.) A seeded random search with six blocks found $84$ further staircase size-passers; none of them falls below that three-block bound (each clears it by at least $0.698$). All of these configurations fail the divisibility conditions, so none of them is a cycle. Size analysis cannot forbid the staircase shape as a cycle for the same reason that analysis of average drift cannot forbid an orbit from growing forever (Remark~\ref{rem:staircase}): in this one configuration, the cycle side and the ever-growing side of the problem meet.

\paragraph{Sharpness (Theorem~\ref{thm:staircase}).} The staircases pass every size test for every number of blocks from $3$ on, and for every size threshold from an explicit threshold number of blocks on.
\begin{itemize}
\item \emph{For every number of blocks from $3$ on} (part (i)). For every $p$ of $16$ or more there is a size-passer with $p$ blocks whose $n$ is at least $1.585\ldots$ to the power $p$ and at most $1.05$ times that, and whose $\gamma$ lies between two fixed numbers, about $3.68$ and $5.14$, the same for every $p$ (Theorem~\ref{thm:staircase}(i) gives them exactly). For $3$ to $15$ blocks, size-passers with $\gamma$ in the same range are found by a computer check of a finite list. For $2$ and $4$ blocks, size-passers are exhibited, with no range claimed; nothing is stated for $1$ block.

Why this makes the fading unavoidable: a trim of this shape says that, for every size-passer, $\gamma$ plus the base-$2$ logarithm of $p$ exceeds $n$ times some factor. Apply it to the staircases of part (i), whose $\gamma$ is at most about $5.14$ while their $n$ is at least $1.585\ldots$ to the power $p$. So, at every $p$ of $16$ or more, the factor must be smaller than ($5.14\ldots$ plus the base-$2$ logarithm of $p$) divided by $1.585\ldots$ to the power $p$. Consequently the exponential loss in the trim is intrinsic to size arguments.
\item \emph{With every number large} (parts (ii)--(iii)). Every odd number below $2$ to the power $71$ is known, by computer verification, to reach $1$ \cite{barina}. A nontrivial cycle never reaches $1$, so all its numbers lie above that range, and the natural test of size arguments puts every number of the cycle above a threshold $X_0$. A configuration's smallest number is tied to $\gamma$. Take the base-$2$ logarithm of that smallest number: it lies within about the base-$2$ logarithm of $p$ of $\gamma$ (Theorem~\ref{thm:staircase}(ii) has the exact statement). So a size-passer whose numbers are all at least $X_0$ needs $\gamma$ at least about the base-$2$ logarithm of $X_0$ minus that of $p$, and $\gamma$ slightly above the base-$2$ logarithm of $X_0$ is enough. For every real $X_0$ of $1$ or more, and every number of blocks from the threshold $p_1(X_0)$ on, there is a size-passer with every number at least $X_0$ whose $\gamma$ exceeds the trim's own requirement ($n$ times the shrinking factor) by an amount between two fixed numbers, about $2.71$ and $4.21$ (Theorem~\ref{thm:staircase}(iii)).

The threshold $p_1(X_0)$ is computed from the best whole-number fractions approximating $1.585\ldots$ (its continued fraction). It never decreases as $X_0$ grows. Sample values: $10$ when $X_0$ is at most $16$; $22$ when $X_0$ is between $2^8$ and $2^{12}$; $36$ at $2^{20}$; $98$ at the threshold Simons and de~Weger used, $5$ times $2^{60}$; $110$ at $2^{71}$; and $309$ at $2^{200}$. (These are samples, not a complete table; the record computes $p_1(X_0)$ at every power of $2$ up to $2^{200}$.) No single formula covering every $X_0$ is claimed, and nothing is claimed below $p_1(X_0)$.
\item \emph{What is attained} (``The threshold, explicitly'', below). At that threshold scale the staircases attain Rozier's real-number optimum (the minimax value of his extremal profile) under every rotation's exact size condition --- that value, and not the inequality of his Lemma~4 (his paper is in French: `Lemme~4'), whose extra term, $2$ to the power $p$, is larger throughout this range than $2$ to the power ($n$ times the shrinking factor). This is the sharpness Simons and de~Weger expected. There the staircases attain the exponential dependence on the number of local minima, and show that the multiplier in front of the exponential in Simons and de~Weger's Lemma~7 cannot be lowered below $0.054$. That multiplier currently depends on $p$; whether it can be made a constant, independent of $p$, is left open, as Simons and de~Weger leave it.
\end{itemize}

\paragraph{What is left.} The staircases pass every size condition, and every one tested fails the divisibility conditions. So ruling out cycles uniformly in the number of blocks requires the divisibility conditions of Proposition~\ref{prop:elim}, not size alone --- the conclusion Halbeisen and Hungerb\"uhler reached when measuring cycles by length \cite{hh}. Every statement about the staircases concerns how far size arguments reach, and none bears on whether cycles exist. The staircases of part~(i) computed in the record all fail the divisibility conditions, as those of part~(iv) do; those of part~(iii) were not tested for divisibility. Parts (i)--(iii) were each checked by independently written code, with exact whole-number arithmetic at every size decision (``Verification and record'', below; Appendix~\ref{app:verify}).

\subsection*{The formal statements}

Fix a hypothetical $p$-block cycle of $F$ with entry depths $m_t \ge 1$, exit valuations $s_t \ge 1$ ($t \in \Z/p\Z$), and set $n = \sum m_t$, $K = \sum s_t + n$, $q = 2^K - 3^n$.

\paragraph{Notation.} Two letters keep two meanings, to match the record or the literature: $M_t$ (Proposition~\ref{prop:elim}) is a partial sum of entry depths, unrelated to the anchor $M(\w)$; and $S_t$ (Proposition~\ref{prop:elim}) and $S_n$ (Hypothesis~\ref{hyp:aeh}) are both partial sums of exponents of $2$, over different ranges. $\Lambda$ is the linear form $K\log 2 - n\log 3$. Symbols that carried two meanings in v3 are renamed in this version; Appendix~\ref{app:concordance} lists them.

\begin{proposition}[elimination identity]\label{prop:elim}
$q > 0$, and for every rotation $r$,
\[ \w_r\, 3^{a_{r-1}}\, q \;=\; R_r \;:=\; \sum_{t=0}^{p-1} 3^{M_t}\, 2^{S_t}\,(2^{s_t}-1), \]
where, in rotation order starting at $r$, $M_t = \sum_{j>t} m_j$ and $S_t = \sum_{j<t}\sigma_j$, with $\sigma_j = s_j + m_{j+1}$, indices cyclic (the step's $\sigma = \vt{C}$, $m_+ = \sigma - s$, of Definition~\ref{def:reduced}; the shift is essential, $m_{j+1}$ and not $m_j$; and $M_t$, a partial sum of entry depths, is unrelated to the anchor $M(\w)$ of Appendix~\ref{app:laws} --- this is the project record's own notation, \texttt{cycles.md} \S12.6.1, kept so that paper and record read together). Hence $q \mid R_r$ and $q \le \min_r R_r$. For the trivial cycle, $R = 4^p - 3^p = q$ identically.
\end{proposition}

This is the classical cycle equation of B\"ohm and Sontacchi \cite{bohm} (see \cite[I, item~26]{lagarias}) in block coordinates; in integer form it reads $2^{m_0}R_0 = 2^K F_n(a) + q$, with $F_n$ Tao's $n$-Syracuse offset \cite[(1.5)]{tao} and $a$ the valuation word of the profile (\texttt{cycles.md} Remark~12.6.1.7).

\begin{proof}
The step identity $2^{\sigma_t}3^{a_t}\w_{t+1} = 3^{d_t}\w_t + (2^{s_t}-1)$ (Definition~\ref{def:reduced}) unrolls $p$ times; imposing closure $\w_p = \w_0$, substituting $d_j = m_j + a_{j-1}$, and cancelling the common power of $3$ yields the display. Positivity: the right side is positive.
\end{proof}

\begin{corollary}[size balance]\label{cor:size}
$2^K = 3^n \prod_t (1 + \chi_t)$ with $\chi_t = (2^{s_t}-1)/(3^{d_t}\w_t)$; if every cycle element exceeds $X$ then $0 < K\log 2 - n\log 3 < p/X$. This is the classical Baker-side engine \cite{eliahou}, obtained in four lines.
\end{corollary}

\begin{proof}
Multiplying the step identity $2^{\sigma_t}3^{a_t}\w_{t+1} = 3^{d_t}\w_t\,(1 + \chi_t)$ around the cycle, the $\w$-products cancel, and $\sum_t \sigma_t = K$ and $\sum_t d_t = n + \sum_t a_t$ leave $2^K = 3^n\prod_t(1+\chi_t)$. Each $\chi_t > 0$, so $K\log 2 - n\log 3 = \sum_t \log(1 + \chi_t)$ lies in $(0, \sum_t \chi_t)$. Finally $3^{d_t}\w_t > A_t = 2^{s_t}x_t$, where $x_t$, the exit of block $t$, is a cycle element, so $\chi_t < 1/x_t < 1/X$. (\texttt{cycles.md} Corollary~12.1.2.)
\end{proof}

\begin{lemma}[ceiling]\label{lem:ceiling}
If $K > \lceil n\LL \rceil$ then $\prod(1+\chi_t) \ge 2$, so some $\chi_t \ge 2^{1/p}-1 \ge (\ln 2)/p$ and some block exit is $< p/\ln 2$. Consequently $K = \lceil n \LL\rceil$ for every nontrivial cycle with $p \le \ln 2 \cdot X_0$, $X_0$ the verified range \cite{barina} --- unconditionally for $p\le5$ by inspecting the odd values $\le 7$.
\end{lemma}

\begin{proof}
$K \ge \lceil n\LL\rceil + 1 > n\LL + 1$ gives $\prod_t(1 + \chi_t) = 2^K/3^n > 2$ (Corollary~\ref{cor:size}), so some $1 + \chi_t \ge 2^{1/p}$, and $\chi_t \ge 2^{1/p} - 1 \ge (\ln 2)/p$; since $\chi_t < 1/x_t$ (proof of Corollary~\ref{cor:size}), the exit $x_t$ of that block is $< p/\ln 2$. As $q > 0$ forces $K \ge \lceil n\LL\rceil$, a nontrivial cycle with $K \ne \lceil n\LL\rceil$ therefore has an element below $p/\ln 2 \le X_0$, which reaches $1$: impossible. For $p \le 5$, $p/\ln 2 < 7.22$, and the odd values $1, 3, 5, 7$ all reach $1$ ($3 \to 5 \to 1$; $7 \to 11 \to 17 \to 13 \to 5 \to 1$), so no nontrivial cycle contains them. (\texttt{cycles.md} Lemma~12.6.2 and Corollary~12.6.3.)
\end{proof}

\begin{theorem}[periods $1$--$3$; rederivations]\label{thm:smallp}
$F$ has no nontrivial cycles of $1$, $2$, or $3$ blocks. These statements are contained in the classical results \cite{steiner,rozier90,sdw,hercher}; the derivations, not the theorems, are the contribution.
\end{theorem}

\begin{proof}[Proof outline]
$p=1$: the fixed-point equation is $\w\,3^a(2^{s+m}-3^m) = 2^s-1$; the sign and size constraints pin $2^s$ into a window of endpoint ratio $(1-3^{-m})/(1-2^{-m}) < 2$, leaving at most one candidate $s$ per $m$; exact search kills $m \le 20{,}000$ and Steiner's theorem \cite{steiner} covers all $m$ via the block--circuit correspondence. $p=2,3$: budget-trim lemmas (casework over rotations' dominant terms, using $\sum s_t = K - n$) force $0 < K\log2 - n\log3 < 2^{6-n/5}$ resp.\ $2^{3-0.115n}$; exact searches to $n = 20{,}000$ find no solutions; any effective irrationality measure for $\LL$ \cite{rhin} excludes the tail. Full details are in the project record: the three classifications at \texttt{cycles.md} \S12.2.3, \S12.5.3 and \S12.7.5, the two trim lemmas at \S12.5.2 and \S12.7.4, and the ceiling lemma that supplies $K = \lceil n\LL\rceil$ at \S12.6.2. The searches are \texttt{experiments/period1\_cycles.py}, \texttt{experiments/period2\_cycles.py} and \texttt{experiments/period3\_cycles.py}.
\end{proof}

\paragraph{The dictionary with $m$-cycles.} In the Terras map $x \mapsto (3x+1)/2$ ($x$ odd), $x \mapsto x/2$ ($x$ even), a block of entry depth $m_t$ and exit valuation $s_t$ is exactly $m_t$ odd terms followed by $s_t$ even terms: by Proposition~\ref{prop:block}(i)--(ii) each of $x_0,\dots,x_{m-1}$ is odd and maps to the next by $(3x+1)/2$, the last of them to the even number $A = 3^d\w - 1$, which halves $s = \vt{A} \ge 1$ times down to the next block entry. So a $p$-block cycle of $F$ has exactly $p$ maximal runs of odd terms, and its local minima are exactly its $p$ block entries, each lying below both its predecessor $2x$ and its successor $(3x+1)/2$ --- including the entries of depth $m_t = 1$, which look like pure falls only in the odd-to-odd orbit. Hence a $p$-block cycle is an $m$-cycle of Simons and de~Weger \cite{sdw} with $m = p$, their $k_i$ our $m_t$, their $\ell_i$ our $s_t$, their $K$ our $n$ and their $K+L$ our $K$; Rozier's number of local minima \cite{rozier91} is the same count. Every $m$-cycle exclusion therefore transfers verbatim: Hercher's theorem \cite{hercher} excludes every nontrivial cycle of $F$ with $p \le 91$ blocks, of which Theorem~\ref{thm:smallp} is a small special case. Such exclusions bound the number of blocks only, and say nothing about how many blocks have depth $m_t \ge 2$ (\texttt{cycles.md} \S12.4).

\paragraph{Configurations.} The proofs below use a cycle only through its block data. A \emph{$p$-block configuration} is a tuple of integers $m_t, s_t \ge 1$ ($t \in \Z/p\Z$) with $n$, $K$, $q$ as above and $q > 0$; its rotation numerators $R_r$ are defined by the display of Proposition~\ref{prop:elim}, and $\gamma := K - \log_2 q$. It is a \emph{size-passer} if $q \le R_r$ for every $r$. Put $x_r := 2^{m_r}R_r/q - 1$ and $x_{\min} := \min_r x_r$. On a cycle, Proposition~\ref{prop:elim} gives $x_r + 1 = 2^{m_r}3^{a_{r-1}}\w_r$, so $x_r$ is the entry of block $r$ and, by the dictionary above, $x_{\min}$ is the least element of the cycle; for a configuration the same unrolling over the rationals makes the $x_r$ the block entries of the rational cycle it defines. Write $\Lambda := K\log 2 - n\log 3$; since $3^n/2^K = 1 - q/2^K = 1 - 2^{-\gamma}$, we have $\Lambda = -\log(1 - 2^{-\gamma}) > 2^{-\gamma}$. Write $\kp := (\LL - 1)/((\LL)^p - 1)$ for the trim's exponent (the record's $\varepsilon_p$) and, for an arc $A$ of consecutive blocks, $w(A) := (\LL - 1)\,m(A) - s(A)$, its $m$-mass at exchange rate $\LL - 1$ minus its $s$-mass. For each rotation $r$ and each position $t$ in rotation order from $r$, let $A_t$ be the arc of the blocks after position $t$ (it ends at block $r-1$, and is empty for $t = p-1$). Then the terms of $R_r$ are
\begin{equation}\label{eq:term}
3^{M_t}\,2^{S_t}\,(2^{s_t}-1) \;=\; 2^{K - m_r + w(A_t)}\,\bigl(1 - 2^{-s_t}\bigr),
\end{equation}
since $M_t = m(A_t)$ and $S_t + s_t = (K - n - s(A_t)) + (n - m_r - M_t)$.

\begin{theorem}[uniform trim; Rozier \cite{rozier91}, Simons--de~Weger \cite{sdw}]\label{thm:uniform}
Every nontrivial $p$-block cycle satisfies, with $\gamma := K - \log_2 q$,
\[ \gamma + \log_2 p \;>\; \frac{(\LL - 1)\, n}{(\LL)^{\,p} - 1}. \]
Consequently the candidate set at every period is finite and explicitly bounded. Combined with Rhin's effective bound \cite{rhin},
\[ K\log 2 - n\log 3 \;>\; \exp\bigl(-13.3\,(0.46057 + \log n)\bigr), \]
every nontrivial $p$-block cycle has $n \le n_0(p)$, the unique solution in $n$ of
\[ \frac{(\LL - 1)\,n}{(\LL)^{\,p} - 1} \;=\; \log_2 p \;+\; \frac{p + 13.3\,(0.46057 + \log n)}{\log 2}, \]
so that $n_0(p) = O(p\,(\LL)^p)$.
\end{theorem}

\begin{proof}
For each rotation, $q \le R_r < p\,2^{\max_t e_t}$, where the exact exponent identity gives $e_t - K = (\LL-1)M_t - m_r - \sum_{j>t}s_j$. The resulting conditions $m_r < \gamma + \log_2 p + \Phi_r$ close into the cyclic recursion $\Phi_{r} = \max(0, \Phi_{r-1} + w_{r-1})$ on arc weights $w_j = (\LL-1)m_j - s_j$, whose cyclic total $(\LL-1)n - \sum s_t$ is negative because $2^K > 3^n$ forces $\sum s_t > (\LL - 1)n$. A max-plus recursion with negative cyclic total must vanish somewhere (otherwise following it once around the cycle would add the total to itself); fix such a reset index $r_0$, so $\Phi_{r_0} = 0$. Walking forward, $\Phi_{r_0+j} \le (\LL-1)\mu_{j-1}$ where $\mu_j$ is the partial sum $m_{r_0} + \cdots + m_{r_0+j}$ (drop the $-s$ terms), so the block bound gives $m_{r_0+j} < (\gamma + \log_2 p) + (\LL-1)\mu_{j-1}$, i.e.\ $\mu_j < \LL\,\mu_{j-1} + (\gamma + \log_2 p)$. Iterating $p$ times from $\mu_{-1} = 0$ yields $n = \mu_{p-1} < (\gamma + \log_2 p)\,\bigl(\LL^{\,p} - 1\bigr)/(\LL - 1)$, which is the display. The argument uses only $q \le R_r$ for every $r$, so the first display holds for every size-passing configuration.
\end{proof}

This is Rozier's Lemme~4 \cite[p.~11, with (5) on p.~10]{rozier91} and Simons and de~Weger's Lemma~7 \cite[p.~58]{sdw} in rotation coordinates, up to an additive constant, with their number of local minima equal to our $p$ by the dictionary above; finiteness at each $p$ is Rozier's Th\'eor\`eme \cite[p.~8]{rozier91}, and also Brox's \cite{brox}. The second display is \cite[Lemma~14 and Theorem~3(d)]{sdw}, whose $K_1(p)/(p\,(\LL)^p)$ decreases to $15.107\ldots$, with the Rhin measure used there \cite[Lemma~12]{sdw}; its derivation from the first display and Rhin's bound, in the present notation, is \texttt{cycles.md} Corollary~12.8.2. We rederive it by another route --- rotation numerators and a max-plus recursion, where Simons and de~Weger chain the minima --- because the recursion identifies the extremal profile that Theorem~\ref{thm:staircase} realizes.

\begin{theorem}[sharpness: the staircase]\label{thm:staircase}
Write $\phi(x) := -\log_2(1 - 2^{-x})$ for $x > 0$, a decreasing involution of $(0,\infty)$.
\begin{enumerate}
\item[(i)] \emph{Every period.} For every $p \ge 16$ there are an integer $n \in [(\LL)^p,\, 1.05\,(\LL)^p]$ and a $p$-block size-passer with $\sum_t m_t = n$, $K = \lceil n\LL\rceil$ and
\[ 3.683012 \;\le\; \gamma \;\le\; 5.140212 . \]
Size-passers in the same bracket exist at every period $3 \le p \le 15$, and size-passers exist at $p \in \{2,4\}$. Consequently, if $\gamma + \log_2 p > c_p\, n$ for every size-passer at a period $p \ge 16$, then $c_p < (5.140212 + \log_2 p)\,(\LL)^{-p}$: the exponential loss in Theorem~\ref{thm:uniform} is intrinsic to size arguments.
\item[(ii)] \emph{Necessity at a size threshold.} Every configuration satisfies
\[ \gamma - 1 \;\le\; \log_2(x_{\min} + 1) \;<\; \gamma + \log_2 p . \]
Hence for real $X_0 \ge 1$, every size-passer with $x_{\min} \ge X_0$ has $\gamma > \max\bigl(\kp n,\ \log_2(X_0+1)\bigr) - \log_2 p$, and every configuration with $\gamma \ge \log_2(X_0 + 1) + 1$ has $x_{\min} \ge X_0$.
\item[(iii)] \emph{Sufficiency at a size threshold.} Let $X_0 \ge 1$ be real, $\lambda_0 := \LL/(\LL-1) + 2^{-20} = 2.7095122\ldots$ and $g := 1 + \max\bigl(\log_2(X_0+1),\, \lambda_0\bigr)$. For every period $p \ge p_1(X_0)$, with $p_1(X_0)$ defined in the proof from the continued fraction of $\LL$, there is a $p$-block size-passer with $x_{\min} \ge X_0$ and
\[ g \le \gamma \le g + 1, \qquad \lambda_0 \;\le\; \gamma - \kp n \;\le\; \lambda_0 + 1.5, \qquad \kp n \;\le\; \rho_0 \;\le\; \gamma , \]
where $\rho_0 := \max_j\bigl(m_j - (\LL-1)(m_0 + \cdots + m_{j-1})\bigr)$ is Rozier's quantity in block order from block $0$. This configuration has $\kp n > g - \lambda_0 - \tfrac12$ and $\Lambda > 2^{-\gamma} > 0.054\cdot 2^{-\kp n}$.
\item[(iv)] \emph{Instances.} At $p=7$, $n=94$, the staircase $m = (4,7,9,15,23,35,1)$ with $s = (1,\dots,1,\,\stot-6)$ passes all seven size tests with $\gamma = 6.74$, where the constant of the period-$3$ trim, $\gamma > 0.1157\,n - 2$ (\texttt{cycles.md} Lemma~12.7.4), would demand $\gamma > 8.9$; a seeded random search at $p=6$ found $84$ further staircase size-passers, none below that constant (each clears it by at least $0.698$); all of these fail the divisibility conditions $q \mid R_r$.
\end{enumerate}
Uniform cycle exclusion therefore requires the divisibility conditions $q \mid R_r$ of Proposition~\ref{prop:elim}, not size alone --- the conclusion Halbeisen and Hungerb\"uhler reached in the cycle-length parameter \cite[p.~238]{hh}.
\end{theorem}

\begin{proof}
\emph{The function $\phi$.} $2^{-x} + 2^{-\phi(x)} = 1$ is symmetric in $x$ and $\phi(x)$, so $\phi$ is an involution; it decreases because $1 - 2^{-x}$ increases.

\emph{Two facts used throughout.} (a) For $K = \lceil n\LL\rceil$ put $\delta_n := K - n\LL \in (0,1)$; then $q/2^K = 1 - 2^{-\delta_n}$, so $\gamma = \phi(\delta_n)$, and $\delta_{n+j} \equiv \delta_n - j\LL \pmod 1$. (b) Multiplying \eqref{eq:term} by $2^{m_r}/q = 2^{m_r + \gamma - K}$ and summing over $t$ gives, for every rotation $r$,
\begin{equation}\label{eq:element}
x_r + 1 \;=\; 2^{\gamma}\sum_{t=0}^{p-1} 2^{w(A_t)}\,\bigl(1 - 2^{-s_t}\bigr).
\end{equation}

\emph{(ii).} In \eqref{eq:element} the empty arc ($t = p-1$) contributes $1 - 2^{-s_{r-1}} \ge \tfrac12$ and every term is positive, so $x_r + 1 \ge 2^{\gamma-1}$ for every $r$. For the upper bound put $W_j := \sum_{i<j}\bigl((\LL-1)m_i - s_i\bigr)$, so $W_0 = 0$ and $W_p = n\LL - K < 0$; choose $j^\ast \in \{1,\dots,p\}$ minimizing $W_j$ and put $r_0 := j^\ast \bmod p$. An arc ending at block $r_0 - 1$ either does not wrap, with weight $W_{j^\ast} - W_i \le 0$ ($0 \le i \le j^\ast$; at $i = 0$ the weight is $W_{j^\ast} \le W_p < 0$), or wraps, with weight $W_p + W_{j^\ast} - W_i < 0$; so every term of \eqref{eq:element} at $r = r_0$ is $< 1$, and $x_{r_0} + 1 < p\,2^{\gamma}$. With Theorem~\ref{thm:uniform}, whose proof uses only the size conditions, a size-passer with $x_{\min} \ge X_0$ has $\gamma + \log_2 p > \kp n$ and $\gamma + \log_2 p > \log_2(x_{\min} + 1) \ge \log_2(X_0 + 1)$; and $\gamma \ge \log_2(X_0+1) + 1$ gives $x_{\min} + 1 \ge 2^{\gamma - 1} \ge X_0 + 1$.

\emph{The construction.} Given $p \ge 2$ and $n \ge 1$, let $K = \lceil n\LL\rceil$, $\stot = K - n$, $q = 2^K - 3^n$, $\gamma = \phi(\delta_n)$, $s_c := \stot - p + 1$ and $\eta := \phi(s_c)$, and assume
\[ \text{(H0)}\quad \stot \ge p \ \text{ and } \ (\LL-1)(n-p) + \gamma \ge p + \eta; \qquad
   \text{(H1)}\quad \gamma \ge \max\bigl(2 + \eta,\ \Gamma'(p,n)\bigr), \]
where
\[ \Gamma'(p,n) := 1 + \eta + \frac{(\LL)^p - (\LL)^2 - (\LL-1)(p-2) + (\LL-1)^2 n}{(\LL-1)\bigl((\LL)^p - 1\bigr)} \;=\; \kp n + \Gamma'_0(p) + \eta, \]
with $\Gamma'_0(p) < \LL/(\LL-1)$, the last inequality being $(\LL)^2 + (\LL-1)(p-2) > 1$. Set $s_t := 1$ for $t \le p-2$ and $s_{p-1} := s_c$, and for $r = 0,\dots,p-2$, with $\mu_{-1} := 0$,
\[ \begin{gathered} f_r := \gamma + (\LL-1)\,\mu_{r-1} - r - \eta, \qquad b_r := n - 1 - \mu_{r-1} - (p-2-r), \\ m_r := \min\bigl(\lfloor f_r \rfloor,\, b_r\bigr), \qquad \mu_r := \mu_{r-1} + m_r, \end{gathered} \]
and $m_{p-1} := n - \mu_{p-2}$, the \emph{crash} block. The floor is an exact integer comparison: $\lfloor f_r\rfloor = \max\{m : q\,2^m \le X_r\}$ with $X_r := 3^{\mu_{r-1}}\,2^{\,n - \mu_{r-1} + p - 1 - r}\,(2^{s_c}-1)$, since $\log_2 X_r = K - r - \eta + (\LL-1)\mu_{r-1}$. We show that every $m_t \ge 1$, that $\sum_t m_t = n$ and $n + \sum_t s_t = K$, that the configuration is a size-passer, and that $\kp n \le \rho_0 \le \gamma - \eta$.

\emph{Well-definedness.} From $p \le \stot < (\LL-1)n + 1$ we get $n > (p-1)/(\LL-1)$, hence $n \ge p$ and $b_0 = n - p + 1 \ge 1$; and $b_{r+1} = b_r - m_r + 1 \ge 1$ since $m_r \le b_r$. So every climb block has $m_r \ge 1$ as soon as $\lfloor f_r\rfloor \ge 1$, and $\mu_{p-2} \le n - 1$, so the crash has $m_{p-1} \ge 1$. While the budget $b_r$ has not bound, $m_r = \lfloor f_r\rfloor$: here $f_0 = \gamma - \eta \ge 2$, and $f_{r+1} - f_r = (\LL-1)m_r - 1 > 0$ whenever $m_r \ge 2$, so $\lfloor f_r\rfloor \ge 2$ throughout. Once the budget binds at some $r$, every later $b_j$ equals $1$, $\mu_{j-1} = n - p + j$, and $f_j = \bigl(\gamma + (\LL-1)(n-p) - \eta\bigr) - (2-\LL)\,j \ge p - (2-\LL)(p-2) \ge 2$ by (H0), so $m_j = 1$.

\emph{The climb rotations $r \le p-2$.} In rotation order from $r$ the crash block sits at position $t = p-1-r$, with $M_t = \mu_{r-1}$ and $S_t = (p-1-r) + (n - \mu_r)$, so its term of $R_r$ is $3^{\mu_{r-1}}2^{\,n - \mu_r + p - 1 - r}(2^{s_c}-1) = X_r/2^{m_r} \ge q$, because $m_r \le \lfloor f_r\rfloor$. Every term is positive, so $q \le R_r$.

\emph{The crash rotation $r = p-1$.} Its position-$0$ term is $3^{\mu_{p-2}}(2^{s_c}-1)$, and this is $\ge q$ if and only if
\[ \text{(D)}\qquad \LL\, \mu_{p-2} \;\ge\; n + p - 1 - \gamma + \eta . \]
If the budget binds somewhere, $\mu_{p-2} = n-1$ and (D) reads $(\LL-1)n \ge \LL + p - 1 - \gamma + \eta$, which follows from (H0) because $(\LL - 1)(p-1) \ge 0$. If it never binds, $m_r = \lfloor f_r\rfloor > f_r - 1$ gives $\mu_r > \LL\,\mu_{r-1} + \gamma - r - \beta$ with $\beta := 1 + \eta$, so by induction $\mu_r > A_r$ for the real recursion $A_r = \LL\,A_{r-1} + \gamma - r - \beta$, $A_{-1} = 0$, whose solution is
\[ A_r \;=\; \frac{(\LL)^{r+1}\bigl((\gamma-\beta)(\LL-1) - 1\bigr) + (r + \beta - \gamma)(\LL-1) + \LL}{(\LL-1)^2}. \]
Clearing denominators, $\LL\,A_{p-2} - (n + p - 1 - \gamma + \eta) = \bigl(((\LL)^p - 1)/(\LL - 1)\bigr)\,\bigl(\gamma - \Gamma'(p,n)\bigr) \ge 0$ by (H1), which gives (D). (The crash depth is $1$ exactly when the budget binds.)

\emph{Rozier's quantity.} The climb terms satisfy $m_r - (\LL-1)\mu_{r-1} \le f_r - (\LL - 1)\mu_{r-1} = \gamma - r - \eta$, and the crash term is $n - \LL\,\mu_{p-2} \le \gamma + 1 - p - \eta$ by (D), so $\rho_0 \le \gamma - \eta$. Conversely, all terms $\le \rho_0$ give $\mu_j \le \LL\,\mu_{j-1} + \rho_0$ and $n \le \rho_0((\LL)^p - 1)/(\LL - 1)$, i.e.\ $\rho_0 \ge \kp n$; this is Rozier's minimax \cite[p.~11]{rozier91}, valid for positive reals in any cyclic order.

\emph{(i).} Let $\vartheta := 8 - 5\LL = 0.0751874964\ldots$, $\Gamma^\ast := 1.05\,\LL(\LL-1) + 1/(\LL-1) + 1 = 3.6830121\ldots$, $\delta_{\mathrm{hi}} := \phi(\Gamma^\ast) = 0.1169390665\ldots$ and $\delta_{\mathrm{lo}} := 0.0415$. Since $5\LL = 8 - \vartheta$, $\delta_{n+5} \equiv \delta_n + \vartheta \pmod 1$. The fourteen points $j\vartheta \bmod 1$, $0 \le j \le 13$, are already sorted ($13\vartheta < 1$), with gaps $\vartheta$ and $1 - 13\vartheta < \vartheta$; so the fourteen closed arcs $[\delta_{\mathrm{lo}}, \delta_{\mathrm{hi}}] - j\vartheta$, of length $0.0754390665\ldots > \vartheta$, have left endpoints at mutual gaps at most $\vartheta$ and cover the circle. Hence for every $N$ one of $\delta_N, \delta_{N+5}, \dots, \delta_{N+65}$ lies in $[\delta_{\mathrm{lo}}, \delta_{\mathrm{hi}}]$: any $66$ consecutive integers contain an $n$ with $\delta_{\mathrm{lo}} \le \delta_n \le \delta_{\mathrm{hi}}$. For $p \ge 16$ the window $[(\LL)^p, 1.05\,(\LL)^p]$ contains at least $0.05\,(\LL)^p - 1 > 78$ consecutive integers, hence such an $n$, and since $\phi$ decreases, $\Gamma^\ast \le \gamma = \phi(\delta_n) \le \phi(\delta_{\mathrm{lo}}) = 5.1402114\ldots$. For this $n$, $\stot > (\LL-1)n \ge (\LL-1)(\LL)^p > p + 900$, so $\eta < 2^{-900}$, and $(\LL-1)(n-p) > p + 1$, which gives (H0); and $\Gamma'(p,n) < 1.05\,(\LL)^p\kp + \LL/(\LL-1) + \eta < 3.33 < \Gamma^\ast \le \gamma$ gives (H1). The construction applies. The periods $3 \le p \le 15$, at exhibited $n$ with $n/(\LL)^p \le 1.70$, and the exhibitions at $p \in \{2, 4\}$, which this argument does not reach, are finite computations, recorded at \texttt{cycles.md} \S12.8.6.1 and \S12.8.6.4. Finally $n \ge (\LL)^p$ and $\gamma \le 5.140212$ give the bound on $c_p$.

\emph{(iii).} Let $P_j/Q_j$ be the convergents of $\LL$, $\Theta_j := |Q_j\LL - P_j|$ and $\Delta := \phi(g) - \phi(g+1)$; let $N_\ast := Q_k + Q_{k-1}$ for the least $k \ge 1$ with $\Theta_{k-1} \le \Delta$, and let $p_1(X_0)$ be the least $p$ with $1/(2\kp) = ((\LL)^p - 1)/(2(\LL - 1)) \ge \max(N_\ast,\, 3p + 40)$ (the condition then holds at every larger $p$). Fix $p \ge p_1(X_0)$, put $n^\ast := \lfloor (g - \lambda_0)/\kp \rfloor$ and take the window $[n^\ast - N_\ast + 1,\, n^\ast]$ of $N_\ast$ consecutive integers. By (a), the points $\delta_n$ on it form a reflected translate of $\{j\LL \bmod 1 : 0 \le j < N_\ast\}$, whose gaps, by the three-distance theorem in its two-length case \cite{sos,swierczkowski}, are $\Theta_{k-1}$ and $\Theta_k$, both at most $\Delta$. So some $n$ in the window has $\delta_n \in [\phi(g+1), \phi(g)]$, i.e.\ $g \le \gamma \le g+1$. On the window $n > (g - \lambda_0)/\kp - N_\ast \ge 1/(2\kp) \ge 3p + 40$, since $g - \lambda_0 \ge 1$; hence $\stot > (\LL - 1)n$ gives $s_c \ge 21$ and $\eta < 2^{-20}$, and $(\LL-1)(n-p) \ge p + 1$, so (H0) holds. Also $\kp n \le \kp n^\ast \le g - \lambda_0$, so $\Gamma'(p,n) < \kp n + \lambda_0 \le g \le \gamma$, and $\gamma \ge g \ge 1 + \lambda_0 > 2 + \eta$: (H1) holds, and the construction gives a size-passer. Since $\gamma \ge g \ge \log_2(X_0+1) + 1$, (ii) gives $x_{\min} \ge X_0$. For the brackets, $\gamma - \kp n \ge g - (g - \lambda_0) = \lambda_0$; and $\kp n > g - \lambda_0 - \kp N_\ast \ge g - \lambda_0 - \tfrac12$, whence $\gamma - \kp n < \lambda_0 + \tfrac32$. Finally $\Lambda > 2^{-\gamma} \ge 2^{-(\lambda_0 + 1.5)}\,2^{-\kp n}$, and $2^{-(\lambda_0+1.5)} = 0.05405\ldots$.

\emph{(iv)} is a computation of the record: Remark~12.8.3 of \texttt{cycles.md}, by the script \texttt{experiments/\allowbreak uniform\_\allowbreak trim.py}. The closing sentence is the content of (i)--(iii): at every period $p \ge 3$, and at every threshold for $p \ge p_1(X_0)$, the size conditions are met at the scale the trim allows, so no argument from size alone excludes cycles uniformly in $p$.
\end{proof}

\paragraph{The threshold, explicitly.} $p_1(X_0)$ is explicit for each $X_0$ through the continued fraction of $\LL$: $p_1 = 10$ for $X_0 \le 16$, $22$ for $2^8 \le X_0 \le 2^{12}$, $36$ at $2^{20}$, $98$ at Simons and de~Weger's $X_0 = 5\cdot 2^{60}$ (consistent with their size-only exclusion of $m \le 75$ there), $110$ at $2^{71}$ and $309$ at $2^{200}$; over $X_0 = 2^j$, $j \le 200$, $p_1 - \log_{\LL} X_0 \in [2.783, 10.630]$ (\texttt{cycles.md} \S12.8.7.4). The values at $X_0 = 2^j$, at $5\cdot 2^{60}$ and at $2^{71}$ are computed by \texttt{experiments/staircase\_x0\_threshold.py}, and the two ranges follow from their endpoints because $p_1$ is nondecreasing in $X_0$ ($\Delta$ decreases as $g$ grows). No closed form uniform in $X_0$ is claimed, and nothing is claimed below $p_1(X_0)$. At the threshold, Simons and de~Weger's Lemma~7 reads $\Lambda < p\,c_p\,2^{-\kp n}$ \cite[p.~58]{sdw}, against the elementary $\Lambda < p/X_0$ of Corollary~\ref{cor:size}; part~(iii) is stated at the scale $\kp n > g - \lambda_0 - \tfrac12$, at which the two cross over (\texttt{cycles.md} Remark~12.8.7.6). There it attains the exponential dependence on the number of minima, and shows that the prefactor $p\,c_p$ cannot be lowered below $0.054$. Whether it can be lowered \emph{to} a constant --- the $\log_2 p$ between parts (ii) and (iii) --- is left open, as Simons and de~Weger leave it \cite[p.~69]{sdw}. The integer profile of (iii) realizes the minimax value of Rozier's optimality remark under every rotation's exact size condition, $\kp n \le \rho_0 \le \kp n + \lambda_0 + 1.5$; what is attained is that value, not the inequality of his Lemme~4, whose additive term $2^p$ exceeds $2^{\kp n}$ throughout this range (\texttt{cycles.md} Remark~12.8.7.6).

\begin{remark}\label{rem:staircase}
The staircase profile --- geometric depth growth with shallow exits --- is exactly the growth regime of divergent orbits. Size analysis cannot forbid it as a cycle for the same reason drift analysis cannot forbid divergence. The two halves of the residual difficulty meet in one configuration.
\end{remark}

\begin{remark}[the crash block]\label{rem:crash}
The crash block of the construction has depth $1$ exactly when the budget binds, and depth at least $2$ otherwise; the record's construction at the canonical scale (\texttt{cycles.md} \S12.8.6.2) assumes a stronger hypothesis that forces depth $1$. At the scale of Theorem~\ref{thm:staircase}(iii) a unit block is excluded wherever the inequality below holds. If a size-passer with $p \ge 2$ has some $m_t = 1$, walking the recursion in the proof of Theorem~\ref{thm:uniform} from its reset, with $\mu_{j_1} = \mu_{j_1-1} + 1$ at the unit block $j_1$ in place of the block bound, gives $n < (\gamma + \log_2 p)\bigl((\LL)^{p-1} - 1\bigr)/(\LL - 1) + (\LL)^{p-1}$, that is,
\[ \gamma + \log_2 p \;>\; \frac{(\LL - 1)\bigl(n - (\LL)^{p-1}\bigr)}{(\LL)^{p-1} - 1} \;\ge\; \LL\,\kp n - (\LL - 1)\,\frac{(\LL)^{p-1}}{(\LL)^{p-1} - 1} : \]
a unit block costs a factor $\LL$ in the slope. Against $\gamma \le \kp n + \lambda_0 + 1.5$ this excludes a unit block wherever $(\LL - 1)\kp n \ge \lambda_0 + 1.5 + (\LL - 1)(\LL)^{p-1}/((\LL)^{p-1} - 1) + \log_2 p$, which at the periods $p_1(X_0)$ and $p_1(X_0) + 20$ holds from $X_0 = 2^{20}$ through $2^{200}$ (\texttt{cycles.md} Remark~12.8.7.5; the table is in the output of \texttt{experiments/staircase\_x0\_threshold.py}); below that the $\log_2 p$ slack leaves the unit crash undecided. Where the condition holds, the crash is forced to be deep: the last block continues the climb, as in Rozier's real extremal.
\end{remark}

\paragraph{Verification and record.}
% PIN-AT-BUILD: the record commit is \recordcommit (preamble), the same as Appendix A.
Prompted by correspondence with Eric Merle concerning the sharpness hedge of this theorem's statement in v1--v3 (his related formal work is \cite{merle}), parts (i)--(iii) were proved in the project record, \texttt{cycles.md} \S12.8.6 and \S12.8.7 at the commit of Appendix~\ref{app:method}, each checked there by independently written code with exact integer arithmetic at every size decision: part~(i) by the scripts \texttt{experiments/staircase\_gamma\_upper.py} ($45$ checks, $0$ failures), \texttt{experiments/\allowbreak staircase\_allp\_construction.py} and \texttt{experiments/\allowbreak staircase\_allp\_diophantine.py} (the sweep, the bracket, the window counts, and the construction at every period from $3$ to $32$ apart from $p = 4$, including the finite range $3 \le p \le 15$; the exhibitions at $p \in \{2,4\}$ are by \texttt{experiments/staircase\_allp.py}, \S12.8.6.4), and parts (ii)--(iii) by \texttt{experiments/staircase\_x0\_threshold.py} ($188$ checks, $0$ failures: the sandwich and the trim on $30{,}000$ random configurations, the two-distance structure, $p_1(X_0)$ for $X_0 = 2^j$, $j \le 200$, and the construction end to end at $95$ pairs $(X_0, p)$), and again by an independent re-implementation, \texttt{experiments/x0\_threshold\_independent\_check.py} ($43$ checks, $0$ failures). The proof printed above composes the record's availability theorem for (i) (\S12.8.6.1) with its construction under the weaker hypothesis of \S12.8.7.2. The configurations of (i) computed in the record all fail the divisibility conditions $q \mid R_r$, as those of (iv) do; those of (iii) were not tested. Every statement here concerns the reach of size arguments, and none bears on whether cycles exist.

\section{The equidistribution hypothesis in full}\label{app:aeh}

This appendix states formally what Section~\ref{sec:aeh} tells in words. It opens with a guided account in plain words; the formal account follows, from ``The window law'' on.

\subsection*{A guided account}

The classical heuristic picture of Collatz orbits (Terras \cite{terras}) treats the halvings as tosses of a fair coin. This account states that picture exactly, as a statement about the shapes of the blocks of typical orbits (the formal account below also states it as a law for finite pieces of the state). It then says which part is proved and which is conjectured, what the statement would and would not give, and how it was tested.

\paragraph{Letters and the model.} Read each block as a \emph{letter}, written $m|s$: the block's entry depth and exit valuation (the bar keeps a letter from being mistaken for a state). In the running example the letters are $3|1$, $1|2$, $1|3$, $1|1$. The model treats the two parts of a letter as independent counts of the kind `toss a fair coin until it lands tails, and count the tosses': each part is $1$ with chance $\tfrac12$, $2$ with chance $\tfrac14$, and so on, halving each time. A letter's chance is the product of its two parts' chances; for example, the letter $1|2$ has chance $\tfrac12 \times \tfrac14 = \tfrac18$. A \emph{word} --- a run of consecutive letters --- has the product of its letters' chances.

\paragraph{What the model says by itself.} On the states, this model becomes an exactly computable law for finite pieces of the state, the \emph{window law} below; it lumps together all values beyond a fixed cut-off, its \emph{cap}. The part of the model that governs factors of $3$ --- the long-run law of the absorption --- is not new here: it is a law Tao introduced, his `Syracuse random variable' \cite{tao}, rescaled to the present setting. Under the window law, unconditionally --- that is, as facts about the model itself:
\begin{itemize}
\item the exit valuation is $1$ with chance $\tfrac12$, $2$ with chance $\tfrac14$, and so on, exactly, at every value below the cap;
\item a factor of $3$ is gained with chance $\tfrac13$, up to the cap: by Section~\ref{sec:anchor} that happens exactly when the exit valuation is even, which has chance $\tfrac14 + \tfrac1{16} + \tfrac1{64} + \dots = \tfrac13$;
\item the classical negative drift holds: on average, the numbers shrink.
\end{itemize}
But these are properties of the model, not of orbits. Matching the model's frequencies, piece by piece at each fixed window size, does not control the averages of entry depths and exit valuations along an orbit: entry depths and exit valuations have no upper limit, and a few very long blocks can dominate an average. So no statement that orbits drift downward or contract follows from it. The usual list of these facts (the folklore `ledger') is thus true of the model; the question is only whether the orbits of typical starting values follow the model.

\paragraph{The hypothesis (Hypothesis~\ref{hyp:aeh}).} Pick a large number and draw an odd starting value at random between it and twice it. Follow its blocks for a number of blocks proportional to its number of binary digits; the proportion is the \emph{horizon rate}. Keep a running total of the halvings done so far, and record each block's letter only while that total stays below a \emph{budget}, also proportional to the number of binary digits (the proportion is the \emph{budget rate}); later blocks are recorded as blanks. For example, for a start with about $40$ binary digits, a horizon rate of $\tfrac1{10}$ follows $4$ blocks, and a budget rate of $\tfrac12$ records letters while fewer than $20$ halvings have been done.

The hypothesis says: for every finite word of letters and every tolerance, the proportion of starts whose recorded frequency of that word differs from the model's by more than the tolerance tends to zero as the starting size grows --- for every \emph{admissible} pair of budget and horizon rates. Admissible means two things.
\begin{itemize}
\item The budget is \emph{protected}: while the budget lasts, the orbit has not yet fallen among the small numbers. Precisely, all but a fraction of starts that shrinks to zero keep every recorded block's exit (the odd number the block ends at) above any level that grows more slowly than every fixed power of the starting size.
\item The pair is \emph{consistent}: under the model each part of a letter averages $2$, so a block uses $4$ halvings on average, and the horizon must fit inside the budget on that average --- four times the horizon rate below the budget rate. (In the example, $4 \times \tfrac1{10}$ is below $\tfrac12$.)
\end{itemize}

\paragraph{What is proved, and what is conjectured.} The hypothesis, for every admissible pair, is a conjecture. Part of it is proved: for every admissible pair whose budget rate is below $1$ --- that is, while the orbit has done fewer halvings than the start has binary digits --- the hypothesis holds, at every word length. This is the \emph{unconditional base case} (``The unconditional base case'', below; proof assembled in the project record, Appendix~\ref{app:verify}). Since the budget rate must exceed four times the horizon rate, a budget rate below $1$ allows exactly the horizon rates below one quarter of a block per binary digit. Beyond that range the hypothesis remains a conjecture.

Why the base case holds: the first blocks of a start are fixed by its own binary digits. The starts whose first blocks follow a given pattern of letters form exactly one class of remainders on division by a power of $2$ (the classical coding fact, Terras \cite{terras}). For a pattern using fewer than ten halvings in all, counting the start's own block, that class is a class of remainders on division by $2$ to a power of at most $10$, and any run of $1024$ consecutive integers contains each such class equally often. So, for example, among the odd numbers of any run of $1024$ consecutive integers, each such pattern occurs in exactly its fair-coin share. The errors from longer patterns, and from the sampling, are explicitly bounded (below); the rest of the proof, which shows that frequencies concentrate, is in the project record and is not reproduced here.

The natural clock here is total halvings, not blocks. The limit of the base case is one halving per binary digit of the start --- the start's own supply of digits --- and that is one quarter of a block per digit only because a block uses $4$ halvings on average. The hypothesis is the assertion that the coin law still describes the orbit after the start's digits have been spent; the digit budget of Section~\ref{sec:anchor} locates that frontier, and the hypothesis is the statistical statement beyond it.

The frontier is the classical one. One quarter of a block per digit is where Terras's counting argument (the classical count of which starts follow which pattern) stops, and Korec's exponent \cite{korec} is what that argument gives when pushed to its limit. It is a limit of this technique, not of the problem. Inselmann \cite{inselmann} goes beyond it without any hypothesis:
\begin{itemize}
\item for the range within which an orbit stays, and for the proportion of its steps that are odd;
\item measured in single steps rather than blocks;
\item using a stronger notion of `almost all', rather than a sharper count.
\end{itemize}

\paragraph{What it would give.} If the hypothesis holds, then for every tolerance and every horizon rate, all but a vanishing proportion of the odd starting values of each size show the model's frequencies along their recorded blocks; all but a vanishing fraction of those blocks lie within the budget, and all of them do when the budget rate is below $1$. Checking each odd start once, at its own size (the power of $2$ just below it), the starts that fail form a set of density zero among the odd numbers --- below any large bound, their share of the odd numbers shrinks to zero (given the hypothesis; ``What the hypothesis implies'', below; proof assembled in the project record, Appendix~\ref{app:verify}). What the hypothesis adds beyond the classical count is the shape of the distribution, and that is the whole of its content: the full halving law of the exit valuation at every value below the cap, the exact $\tfrac13$ rate of factor-$3$ gains, and the law of the depths --- at every finite depth, and beyond the budget at which the classical count stops.

\paragraph{What it would not give.}
\begin{itemize}
\item Control of all sizes at once: the starts that fail at some size or other, taken over every sampling size together, form a larger set, which the hypothesis does not control.
\item Anything about whole orbits: the starts it allows to misbehave are a vanishing proportion of starting values at a fixed, finite horizon --- not a negligible set of orbits --- and it says nothing about where any single orbit ends up in the long run.
\item The exclusion of any single orbit shaped like a staircase.
\item Repeated use along an orbit: the numbers reached from a set containing almost all starts need not contain almost all numbers at the next size (an obstruction Tao states in the same terms \cite{tao}).
\item A proof by finite-window computation: by Heuristic~\ref{prop:budget}, the hypothesis is not expected to be reachable that way.
\end{itemize}
Nor does the paper state the descent of almost every orbit as a consequence of the hypothesis: that descent is already a theorem without it, and a stronger one (Inselmann \cite{inselmann}). The hypothesis's real content is a question about the digits of the anchors of Section~\ref{sec:anchor} --- in the usual terms, digits of $2$-adic logarithms of whole numbers (Appendix~\ref{app:laws} names them precisely) --- and questions of that kind are older and broader than the Collatz problem.

\paragraph{How it was tested.} An eight-round test campaign looked for departures from the model in real orbits. It found four apparent ones, and each dissolved under a sharper control:
\begin{enumerate}
\item[(1)] an excess of exit valuation $3$, and
\item[(2)] correlations between neighbouring blocks --- both artifacts of including small starting values;
\item[(3)] biases in the binary digits of up to $41$ standard deviations (a standard deviation measures how far a count strays from what chance predicts), from orbits that had reached the `bottom regime' of small numbers, where particular numbers rather than a statistical law dominate;
\item[(4)] a bias of five standard deviations in one class of states, which survived discarding every number below $2$ to the power $40$. It dissolved when every orbit was followed for the same number of blocks and every visit counted equally ($0.2503 \pm 0.0015$ measured, against $0.25$ expected, over $84{,}739$ visits). It was explained exactly by a fixed routing of remainders between kinds of state (Lemma~\ref{lem:routing}), which the earlier method --- averaging, orbit by orbit, ratios of counts --- had picked up as a bias.
\end{enumerate}
No remaining departure was found at any tested depth or class of states, under the stated protocol. The campaign has three limits. It tests words of one and two letters only. Its deciding runs pool visits across orbits, so they measure the average over orbits of each start's frequencies, not each start's frequencies themselves. And the filter it used to keep orbits away from small numbers looks at the next state's odd core rather than at the block itself, and does remove some visits at the sizes tested (``Calibration'', below).

\subsection*{The window law}

What is the right comparison object? By Theorem~\ref{thm:onestep} every dynamical decision at horizon one is a function of the depth-$k$ window, so the null model is a measure on windows, and the structure forces the choice: the $\w$-residues carry the anchor digits, for which Haar measure is the canonical reference (this is where the classical $2$-adic shift-conjugacy heuristics live), while the depth is \emph{dynamical} --- it is produced by Theorem~\ref{thm:depth} --- and receives instead its own exact renewal law, $\dnext = m_+ + a_+$ with $m_+$ geometric$(1/2)$ and \emph{independent} of the absorption $a_+$, so that it is the distribution of the depth, and not that of the absorption, which is the convolution of the two. The stationary $3$-adic law governing $a_+$ (\texttt{aeh.md} \S13.6.5) is not new here: it is Tao's Syracuse random variable $\mathrm{Syrac}(\Z_3)$ \cite[Lemma~1.12 and Remark~1.13]{tao}, in the present normalisation. Precisely, the $3$-adic past-limit of the block coordinate has the law of $\mathrm{Syrac}(\Z_3)/2$ --- the two differ by the $3$-adic unit $2^{-1}$ because one block here is a run of Syracuse steps terminated by the first exponent $\ge 2$, so that the block chain sees the Syracuse chain conditioned on that event --- and hence $a_+ = \vth{\mathrm{Syrac}(\Z_3) + 2}$, with $P(a_+ = 0) = \tfrac23$, $P(a_+ = 1) = \tfrac{19}{63}$ and $P(a_+ \ge 2) = \tfrac{2}{63}$ read off the nine values of $\mathrm{Syrac}(\Z/9\Z)$ that Tao computes. What the present coordinates add is the \emph{joint labelled} law (\texttt{aeh.md} \S13.6.3(v)): the $\w$-residues are Haar-uniform among odd residues and independent of the depth, the whole carried on the residues \emph{together with} the stratum labels $(s,\sigma,a_+)$ --- $a_+$ being a $3$-adic function of $\w$ (Lemma~\ref{lem:absorption}) which the finite $2$-adic residue window does not determine. Two integers fix the observable: a depth $k$ and a \emph{cap} $D$, chosen together and quantified over. The \emph{capped depth-$k$ window} at a visit is
\[
  W_{k,D} \;=\; \bigl(\, \w \bmod 2^{k+2},\; d \wedge D,\; s \wedge D,\;
  \sigma \wedge D,\; a_+ \wedge D \,\bigr), \qquad u \wedge D := \min(u,D),
\]
a finite alphabet of at most $2^{k+1}D^3(D+1)$ letters; the cap is what bounds it, and it must cap the labels as well as the depth, since $\sigma$ and $a_+$ are unbounded. It is not the window of Theorem~\ref{thm:onestep}, whose alphabet is countably infinite and which decides the next step: $W_{k,D}$ decides nothing, and no statement in this appendix asks it to. Let $\pi_{k,D}$ denote the law of $W_{k,D}$ under the two-sided Bernoulli measure $\hat B = \bigotimes_{i \in \Z}(\mathrm{geom}(\tfrac12) \times \mathrm{geom}(\tfrac12))$ on the door-letter alphabet --- two-sided because the absorption $a_+$ is a function of the orbit's $3$-adic past and has no realisation on the one-sided $2$-adic side (\texttt{aeh.md} \S13.6.3(iii), \S13.6.5); reading the stationary Syracuse variable as the outcome of an iteration extending to arbitrarily large negative times is Tao's own alternative to his positive-time indexing \cite[Remark~1.13, footnote~4 of arXiv~v7]{tao}. ``Product'' names exactly two clauses of $\pi_{k,D}$ --- residue $\perp$ depth, and $m_+ \perp a_+$ --- and no others: $\pi_{k,D}$ is not a product across its coordinates ($s$ is an exact function of the full state), and the process $(W_{k,D}(n))_n$ is stationary but not independent across time (e.g.\ from $(\w \equiv 1\ (8),\ d$ odd$)$ the next depth is exactly $1$). Under $\pi_{k,D}$, unconditionally: $P(s = j) = 2^{-j}$ at every $j < D$, the mass beyond the cap sitting in one tail cell; the $3$-gain rate is $\sum_{j\,\mathrm{even}} 2^{-j} = \tfrac13$ by Lemma~\ref{lem:absorption}, exact up to that same cell, whose parity split the cap hides; and the classical negative drift is a property of $\pi_{k,D}$ --- but only of $\pi_{k,D}$: window equidistribution at each fixed $(k,D)$ does not control the means of the unbounded $m_+$ and $s$, so no drift or contraction statement about orbits follows from it (\texttt{aeh.md} \S13.3.2). The folklore ``ledger'' is thus a theorem \emph{about} $\pi_{k,D}$; the empirical question is only whether the orbits of typical starting values follow $\pi_{k,D}$.

\subsection*{The hypothesis}

\begin{hypothesis}[AEH, ensemble form]\label{hyp:aeh}
Fix a \emph{budget rate} $\tau > 0$ and a \emph{block horizon rate} $\theta > 0$;
for each $N$ put $b = \lfloor\log_2 N\rfloor$, $E_N = \lceil \tau b\rceil$
and $H_N = \lceil \theta b\rceil$. Draw $x$ uniformly from the odd integers of
$[N, 2N)$; put $(\w_0,d_0) = R(x)$ and $(\w_{n+1},d_{n+1}) = F(\w_n,d_n)$; let the
\emph{letter} at block $n$ be $\ell_n = (m_{+,n},\,s_{n+1})$, and let
$S_n = \sum_{i<n}(m_i + s_i)$ be the \emph{exponent spent} before block $n$ ---
the number of $2$'s divided out from $x$ to $x_{\mathrm{exit}}(n-1)$. Tally block
$n$ at the symbol
\[
  \tilde\ell_n \;=\; \ell_n \ \text{ if } S_n < E_N, \qquad
  \tilde\ell_n \;=\; \dagger \ \text{ otherwise.}
\]
For a finite word $w = (w_1,\dots,w_L)$ of letters, let $f_N(w,x)$ be the
frequency of $w$ among the $H_N - L + 1$ blocks
$(\tilde\ell_n,\dots,\tilde\ell_{n+L-1})$, $0 \le n \le H_N - L$: every
block counted exactly once, at weight $1/(H_N-L+1)$ --- the same number for
every block of every orbit, so no block is reweighted by the orbit it came from
and no denominator is random. Then for every finite word $w$ and every
$\varepsilon > 0$,
\[
  \frac{2}{N}\,\#\bigl\{\, x \text{ odd},\ N \le x < 2N \;:\;
  \bigl| f_N(w,x) - B[w] \bigr| > \varepsilon \,\bigr\}
  \;\longrightarrow\; 0 \qquad (N \to \infty),
\]
where, writing $w_i = (m_i, r_i)$ for the components of the $i$-th letter of $w$,
$B[w] = \prod_{i} 2^{-(m_i + r_i)}$, and $B[w] = 0$ for any $w$ containing
$\dagger$, for every admissible pair $(\tau,\theta)$ in the sense defined below.
\end{hypothesis}

Equivalently, and this is the form the calibration measures: for every $k$, $D$
and $L$, the empirical distribution of the $L$-blocks of consecutive capped
windows $W_{k,D}$ over the same tallied blocks, with $\dagger$ carried through
and given mass $0$ by $\pi^{(L)}_{k,D}$, converges in total variation on the
finite window alphabet, off a vanishing density of starts, to
$\pi^{(L)}_{k,D}$, the $L$-block law of the stationary process under $\hat B$.
The equivalence is Theorem 13.6.4 of \texttt{aeh.md}, a deterministic dictionary
between letters and labelled window blocks; the case $L = 1$ recovers
$\pi_{k,D}$, and $\pi^{(L)}_{k,D} \neq \pi_{k,D}^{\otimes L}$. A sampled segment
has no infinite past at its first block, so the reconstruction of
$W_{k,D}$ from letters is exact only away from the start; the
deviation occupies $O(1)$ blocks of a horizon $H_N \to \infty$ and contributes
$O(1/H_N)$ to every frequency above.

There is one limit here, $N \to \infty$, and the sample grows because the sampling
scale grows rather than because any one orbit is run forever. That is forced: along
a fixed convergent orbit the unrestricted empirical distribution is false, the
tail sitting at $(1,1)$ forever, while above a fixed cut such an orbit
supplies only finitely many qualifying visits and none at all once the cut exceeds
its maximum, so a limit in orbit length is empty rather than merely delicate, in
whichever order it is taken (\texttt{aeh.md} \S13.6.6).

\subsection*{The budget and admissibility}

The horizon does the job a bulk cut would do, and does it without one. A step of
the one-division map $y \mapsto y/2$ or $(3y+1)/2$ never lowers $\log_2$ by more
than $1$, so $\log_2 x_{\mathrm{exit}}(n-1) \ge \log_2 x - S_n$ for every start and
every $n$, with no exceptional set; inside a budget $\tau < 1$ of the start's own
bits every tallied exit therefore exceeds $N^{1-\tau}$, far above the
\emph{bottom regime} of \texttt{aeh.md} \S13.1, the fixed drainage basin of small
integers. That matters because any altitude
threshold is a selection on the observable itself --- a cut on
$x_{\mathrm{exit}} = (3^{d}\w - 1)/2^{s}$ censors large $s$, and a cut on the core
$\wnext$ censors large $s$, $m_+$ and $a_+$ together --- whereas the budget clause
$S_n < E_N$ is \emph{predictable}, decided by blocks strictly earlier than
the one it admits, and the tally denominator is the deterministic $H_N$ that
\texttt{aeh.md} \S13.5's standing rule was written to secure. Call $\tau$
\emph{protected} when all but a
vanishing density of starts keep every in-budget exit above any $X_N$ with
$\log X_N = o(\log N)$, and $(\tau,\theta)$ \emph{consistent} when
$\theta\,\mathbb{E}_B[m+r] < \tau$; \emph{admissible} means both. Every
$\tau < 1$ is protected outright by the bound above, and every
$\tau < (1-\log_2\sqrt3)^{-1} = 4.8188\ldots$ is protected at natural density one
by Inselmann \cite[Thm.~1.1]{inselmann}, whose two-sided envelope is stated for
the one-division map --- the same unit, nothing converted.
Consistency is compatibility with the target law rather than a claim about
orbits: under $B$ a block spends $\mathbb{E}_B[m+r] = 4$ of exponent, so
$4\theta < \tau$ says exactly that $B$ itself predicts no out-of-budget block,
which is what $\pi_{k,D}$'s giving $\dagger$ no mass requires. Where the cylinder
count runs it is more, and unconditionally: for $\tau < 1$ no block of the horizon
leaves the budget and the empirical exponent mean converges to $4$
(\texttt{aeh.md} Lemma 13.2.4(g), whose two error terms are, on that range,
precisely the two clauses of admissibility). Past that range the hypothesis
supplies less, and only about frequencies: a vanishing frequency of $\dagger$ says
all but $o(H_N)$ of the first $H_N$ blocks are within budget, and is not a bound on
a sum over its complement --- the block at which the budget is exhausted is tallied
at full weight and its letter is unbounded, and the $o(H_N)$ blocks past it are
uncontrolled, so $H_N^{-1}\sum_{n<H_N}(m_n+r_n) \to 4$ does not follow there and
can fail by any amount: a single letter of size $cH_N$ at the exhausting block, with
the $o(H_N)$ blocks after it tallied at $\dagger$, leaves every fixed pattern
frequency at its $B$-value and moves the mean from $4$ to $4 + c$ (\texttt{aeh.md}
\S13.2). Hence $4\theta < \tau < 1$ is $\theta < 1/4$ and
$4\theta < \tau < 4.8188\ldots$ is $\theta < 1/\beta = 1.2047\ldots$, where
$\beta = 2(2-\LL) = 0.8301\ldots$ is the classical per-block contraction rate and
$4/\beta = (1-\log_2\sqrt3)^{-1}$ an identity --- arithmetic on the definition of
consistency, with the divisor in either block reading the mean of the target law:
a theorem about orbits below the budget, and not a consequence of the hypothesis
above it.

\subsection*{The unconditional base case}

The hypothesis has an unconditional base case, and it is Heuristic~\ref{prop:budget}
with a number on it. The classical coding fact --- the odd integers whose first $n$
blocks realize a prescribed itinerary form exactly one residue class modulo
$2^{S+1}$, $S$ the itinerary's total exponent (Terras \cite{terras}; in the present
coordinates \texttt{itinerary.md} \S14.15.1.5) --- makes the first $n$ blocks of a
start drawn uniformly from \emph{any} window of $2^{J}$ consecutive integers
exactly product-distributed on the event $S + 1 \le J$: the event that the itinerary
has not outspent the start's supply of binary digits. That fact is about the word
beginning at $x$ itself, while $\ell_0$ is the letter of the block \emph{after} the
start's own; it is applied here to the \emph{extended} $(n+1)$-letter word --- the
start's own letter followed by $\ell_0,\dots,\ell_{n-1}$, of total exponent exactly
$S_{n+1}$, the divisions from $x$ down to $x_{\mathrm{exit}}(n)$ --- of whose law
the one displayed is a marginal. Writing
$b = \lfloor\log_2 N\rfloor$, this gives, for every $J \ge 2$,
\[
  \begin{gathered}
    \bigl\lVert \mathrm{Law}(\ell_0,\dots,\ell_{n-1}) - B^{\otimes n}
      \bigr\rVert_{\mathrm{TV}}
    \;\le\; \frac{2^{J+2}}{N} + P_B(S_{n+1} \ge J), \\[2pt]
    P_B(S_{n+1} \ge J) = P\bigl(\mathrm{Bin}(J-1,\tfrac12) < 2(n+1)\bigr),
  \end{gathered}
\]
the second identity because $S_{n+1}$ is the waiting time for the $2(n+1)$-th
head in a fair coin sequence. The first term is the price of a general window
$[N,2N)$ over a dyadic one and is negligible for $J \le (1-\zeta)b$, $\zeta > 0$ fixed; the second
decays at the exact rate $\log 2 - \Ent(2\theta)$ per bit at $n = \lceil\theta b\rceil$,
$\Ent$ the binary entropy --- positive for every $\theta < 1/4$ and \emph{zero} at
$\theta = 1/4$ (\texttt{aeh.md} \S13.2). The bound is on the joint law of the whole
length-$n$ word, so it controls every finite block length at once, and with the
concentration of empirical pattern frequencies, the removal of the segment's initial
past-boundary and the altitude bound above --- which makes the budget its own
protection --- it yields the hypothesis outright:
Hypothesis~\ref{hyp:aeh} is a \emph{theorem} for every admissible $(\tau,\theta)$
with $\tau < 1$, equivalently for every horizon rate $\theta < 1/4$, at every
block length. The assembly is Lemma 13.2.4 of \texttt{aeh.md}.

The natural clock here is total exponent, not blocks: the frontier is one unit of
exponent per bit of start --- the start's own digit supply, with no estimate in it
--- and it reads as $1/4$ blocks per bit only through the mean exponent per block,
$\mathbb{E}_B[m+r] = 4$, a conversion available here precisely because the word is
exactly $B$ here. Hypothesis~\ref{hyp:aeh} is the assertion that the law still
describes the orbit after the start's digits have been spent; the digit budget
locates the frontier, and this is the statistical statement on the far side of it.
That frontier is the classical one --- the same $\theta < 1/4$ is where Terras's
count stops, and $1 - \beta/4 = \log_4 3$ is exactly the exponent Korec
\cite{korec} obtains by exhausting it --- and it is a frontier of this technique
and not of the problem, Inselmann \cite{inselmann} crossing it unconditionally for
the trajectory envelope and the first moment, in step time and by a stronger
density notion rather than a sharper count (\texttt{aeh.md} \S13.3.2).

\subsection*{What the hypothesis implies, and what it does not}

AEH implies the ledger, in the form the hypothesis has: $s$ is a coordinate of
the observable, so $P(s = j) = 2^{-j}$ is a marginal of $\pi_{k,D}$ rather than a
deduction, exactly at every $j < D$ and up to the cap's single tail cell above;
and for every $\varepsilon$ and every horizon rate, all but a vanishing density of
the odd starting values of each size carry those frequencies along their first
$H_N$ blocks, all but $o(H_N)$ of them within the budget --- and all of them when
$\tau < 1$. Evaluating
each odd $x$ once, at its own dyadic scale $2^{\lfloor \log_2 x\rfloor}$, turns
that array of statements into a single one: the odd integers that fail at their
own scale have natural density zero (\texttt{aeh.md} Proposition 13.2.5). The
union over \emph{all} sampling scales is a different and larger object and is not
controlled --- the bad sets form a triangular array, each testing a horizon and a
budget that move with the scale, and vanishing density in every window does not
thin their unrestricted union. The
descent consequence is not stated here, because it is a theorem without the
hypothesis and a stronger one: Inselmann's \cite[Cor.~1.4, Thm.~1.10]{inselmann},
cited in Appendix~\ref{app:related}, are unconditional and run to $4.8188\ldots$ times the
classical range at which this appendix's cylinder count stops, the two measured in
the same units, the second two-sided and uniform in the time. His
\cite[Thm.~1.6]{inselmann}, a density statement about odd steps, is what
carries his own passage from one-division-map time to Syracuse time; the further passage to
block time needs the frequency with which a Syracuse step ends a block, a
two-letter statistic of the parity word that neither theorem supplies and that
$\pi_{k,D}$ does, so it is not available to underwrite the hypothesis: inside the
digit budget it is a theorem of the cylinder count (\texttt{aeh.md}
Lemma 13.2.4(g)), and past it neither a theorem nor a consequence of the hypothesis
(\texttt{aeh.md} \S13.2.3). What the hypothesis supplies beyond these is
distributional and is the whole of its content: the full $2^{-j}$ marginal at every
$j$ below the cap, the exact $\tfrac13$ rate of Lemma~\ref{lem:absorption}, and the
depth law, at every finite depth and past the budget at which the classical count
stops. Its exceptional set is a set of \emph{starting values} of vanishing density
at a prescribed finite horizon --- not a null set of orbits, and not a statement
about any orbit's infinite tail. It does not exclude individual staircase tails
(Remark~\ref{rem:staircase}); it does not iterate --- the image of a density-one set
of starts need not be density-one at the next scale, an obstruction Tao states in the
same terms \cite{tao}; and, by Heuristic~\ref{prop:budget}, it is not expected to be
reachable by finite-window computation. Its content is a question about digits of $2$-adic
logarithms, older and broader than the Collatz problem.

\subsection*{Calibration}

An eight-round campaign (per-orbit statistics over each cell's own conditional counts, $1{,}600$--$2{,}600$ independent orbits per cell, under an altitude cut on the core) produced four apparent violations of $\pi_{k,D}$-uniformity, each dissolved by a sharper control: (1) an $s{=}3$ marginal excess and (2) pair correlations, both artifacts of small starting values; (3) digit biases, contamination by the bottom regime, where deviations reach $z=41$ but reflect particular numbers rather than a measure; and (4) a $5\sigma$ digit bias in one cell surviving a $2^{40}$ cut (code: \texttt{experiments/aeh\_calibration.py}, and for the fourth \texttt{experiments/aeh\_anomaly.py}). The last dissolved under fixed-horizon unweighted sampling ($0.2503\pm0.0015$ against null $0.25$ over $84{,}739$ visits) and was explained exactly:

\begin{lemma}[deterministic routing]\label{lem:routing}
For $\w \equiv 1 \pmod 8$, $d = 1$: here $s = 1$, $C = 3\w + 1$, $\vt{C} = 2$, $a_+ = 0$, so $\wnext = (3\w+1)/4$ and $\dnext = 1$; the residue $\w \bmod 32$ determines the next class exactly: $1 \mapsto [1,1]$ (immediate return to the same class), $9 \mapsto [7,1]$, $17 \mapsto [5,1]$, $25 \mapsto [3,1]$, the lifting class. Here $[r,1]$ denotes the class $\w \equiv r \pmod 8$, $d = 1$, not a state.
\end{lemma}

Ratio estimators over correlated visit sequences inherit this routing (residue $25$ leads into deep descent, hence fewer qualifying visits, hence inflated per-orbit weight), manufacturing a bias aimed at one residue of the class $[1,1]$, whose residue $1$ returns to it deterministically (a self-loop channel). No residual discrepancy was detected at any tested depth or cell, under the stated protocol and within three limits the campaign does not reach past: it tests block lengths $L = 1$ and $L = 2$ only and is silent above; its adjudicating runs pool per visit across orbits, so they measure the across-orbit average of the per-start frequencies --- a consequence of Hypothesis~\ref{hyp:aeh} rather than the per-start statement itself, which would be tested by the distribution across orbits of $\lVert\nu - \pi_{k,D}\rVert_{\mathrm{TV}}$ and is not among the runs reported; and its altitude guard cuts on the core $\wnext$ rather than the door, censoring $s$, $m_+$ and $a_+$ at once, and does bind at finite size, removing $2.6\%$ of visits in the adjudicating run and moving the reported statistics at the third decimal (\texttt{aeh.md} \S13.4, \S13.5). Hypothesis~\ref{hyp:aeh} carries no cut of any kind; it tallies within the exponent budget, which bounds every tallied exit from below without selecting on the letter being tallied.

\section{Related work in detail}\label{app:related}

Block-like decompositions and stochastic models are classical (Terras \cite{terras}; see Lagarias \cite{lagarias} for the literature), and the $2$-adic conjugacy of the Collatz map to the shift underlies all Haar-genericity heuristics. The unconditional density line descending from Terras is long and is not what Appendix~\ref{app:aeh} adds to. Terras \cite{terras} and Everett \cite{everett} place descent below the start at natural density one; Korec \cite{korec} lowers the target to $y^{c}$ for every $c > \log_4 3 = 0.7924\ldots$, in natural density and within $O(\log y)$ steps; Inselmann \cite{inselmann} carries a two-sided trajectory law --- the orbit tracked to within $m^{\pm\varepsilon}$ of its classical drift, simultaneously at every time up to $(\log_2\tfrac43)^{-1}\log_2 m$ Syracuse steps, the full descent horizon in that unit --- at natural density one, and so reaches $m^{\varepsilon}$ for every fixed $\varepsilon>0$; and Tao \cite{tao} reaches every divergent threshold, in logarithmic density. All of these horizons are counted in steps, and their reading in the reduced blocks of Appendix~\ref{app:aeh} is not free: it divides by a mean number of steps per block which is a theorem of the cylinder count inside the digit budget and, past it, neither a theorem nor a consequence of Hypothesis~\ref{hyp:aeh}. What Hypothesis~\ref{hyp:aeh} asserts is not descent but the letter statistics themselves --- equivalently the exact window law $\pi_{k,D}$ --- at horizons past the digit budget of Heuristic~\ref{prop:budget}. Tao's paper carries a footnote to two further $3$-adic studies, and neither is the object appearing here. Wirsching \cite{wirsching} completes the root of the predecessor tree $3$-adically and studies the inverse map as a Markov chain on $[0,1] \times \Z_3^{\times}$ in order to \emph{count} predecessors; his inverse branches therefore carry equal weight and his limiting $3$-adic law is Haar, whereas the law appearing in Appendix~\ref{app:aeh} weights each branch by the probability $2^{-a}$ that an orbit takes it and is consequently not uniform --- $(\tfrac23,\tfrac13)$ at level one where Haar gives $(\tfrac12,\tfrac12)$. The two programmes are nonetheless siblings in shape: what Wirsching needs is that a Collatz-generated object be unbiased in $3$-adic coordinates, and Hypothesis~\ref{hyp:aeh} is that one is unbiased in $2$-adic coordinates. Thomas \cite{thomas} proves that the integers whose orbit reaches $1$ with all six frequencies over the units modulo $18$ below $\tfrac16 + 0.0215$ have counting function at most $N^{0.9999}$; he constructs no limiting measure. The stationary $3$-adic law of Appendix~\ref{app:aeh} (Tao's $\mathrm{Syrac}$) is far from the near-uniformity he rules out: at level one it puts mass $\tfrac23$ on the residue $2 \bmod 3$, which three of his six classes share, so a sequence with those frequencies has some class at frequency at least $\tfrac29 > \tfrac16 + 0.0215$. That his counting bound follows from the law is not claimed. Two-adic and Diophantine methods dominate the cycle literature \cite{steiner,rozier90,sdw,hercher,eliahou}; effective linear forms in logarithms \cite{yu,bl,rhin} are the engine, and the verified range of the conjecture stands at $2^{71}$ \cite{barina}. The linear cycle equation behind Proposition~\ref{prop:elim} is B\"ohm and Sontacchi's \cite{bohm}.

The trim of Theorem~\ref{thm:uniform} is classical. For every number of local minima of a cycle --- which is our number of blocks $p$ (Appendix~\ref{app:cycles}) --- Rozier \cite[Lemme~4]{rozier91} bounded the number of odd steps, with the exponent $(\LL-1)/((\LL)^p-1)$; Rozier \cite[Th\'eor\`eme, p.~8]{rozier91} and Brox \cite{brox} proved finiteness for each $p$; and Simons and de~Weger \cite[Lemmas~7, 14 and Theorem~3]{sdw} made the bound explicit and computed with it. Our proof is a rederivation by another route. Rozier observed that his bound is optimal over real run-length tuples, the extremal tuple growing geometrically at ratio $\LL$ \cite[p.~11]{rozier91}; Simons and de~Weger expected the exponential dependence on the number of minima to be unavoidable \cite[p.~69]{sdw}. In the cycle-length parameter, Halbeisen and Hungerb\"uhler \cite{hh} proved the size criterion optimal, attained by explicit rational cycles of every length, and reduced the remaining question to divisibility \cite[p.~238]{hh}. Theorem~\ref{thm:staircase} gives integer configurations meeting every rotation's exact size condition at every period $p \ge 3$, and, for $p \ge p_1(X_0)$, with every element at least a threshold $X_0$; at that threshold scale they realize the minimax value of Rozier's extremal --- the sharpness Simons and de~Weger expected.

Two contemporary human--AI efforts are close in spirit: a documented exploration \cite{llmcollatz} develops a burst--gap decomposition and an orbit equidistribution conjecture, and a Lean~4 formalization \cite{merle} proves cycle exclusion conditionally on Baker-type and verification hypotheses, including an impossibility lemma for measure-based uniform bounds kindred to our Theorem~\ref{thm:staircase}; neither contains the anchor machinery or the exact per-step laws. A bachelor's thesis of Rackl \cite{rackl} (Klagenfurt, 2021) applies valuation methods to Syracuse cycles via the classical cycle equation; it likewise contains none of the present apparatus.

\section{Verification index}\label{app:verify}

Every computational claim of the paper, with the check that supports it as stated where the claim is made. ``Record'' is a section of the project record at the commit of Appendix~\ref{app:method}; scripts are in its \texttt{experiments/} directory and run as \texttt{python experiments/<name>.py}.

{\small
\begin{longtable}{>{\raggedright\arraybackslash}p{0.2\textwidth}>{\raggedright\arraybackslash}p{0.4\textwidth}>{\raggedright\arraybackslash}p{0.32\textwidth}}
\hline
Statement & What was checked & Script; record \\
\hline
\endhead
Theorem~\ref{thm:vlaw}, Lemma~\ref{lem:absorption}, Theorem~\ref{thm:depth} (Remark~\ref{rem:verify1}) & Independently written code, zero discrepancies: the valuation law on $8{,}000$ random states across both lifting components; the entry-depth law exhaustively on all $63{,}000$ valid states with $\w < 3000$, $1 \le d < 64$ ($\w = 1$ included, where every law holds as stated); the absorption law on $327{,}980$ states ($983{,}954$ comparisons), with the branches $d = h(s)$ and $d < h(s)$ constructed rather than sampled & \texttt{depth\_law\_reverify.py}, \texttt{absorption\_law.py}; \texttt{stage3.md} \S11.8.6.3 and its status header \\
Theorem~\ref{thm:deltaM} (Remark~\ref{rem:verify2}) & The low-order anchor-increment law & \texttt{anchor\_increment.py} \\
Theorem~\ref{thm:onestep} (Remark~\ref{rem:verify2}) & Along real orbits, $k \in \{4,8\}$: $21{,}296$ steps, $21{,}000$ decided with zero errors, $296$ flagged deep with zero violations of $s_+ \ge k+2$ & \texttt{one\_step\_propagation.py} \\
Undecided rate (after Theorem~\ref{thm:onestep}) & Measured $0.0275$ against $0.03125$ at $k = 4$, $0.0019$ against $0.00195$ at $k = 8$ & \texttt{stage4.md}, Remark ``undecided rate'' \\
Lemma~\ref{lem:ceiling}, $p \le 5$ & The odd values $1, 3, 5, 7$ reach $1$; the orbits are printed in the proof & by inspection \\
Theorem~\ref{thm:smallp} & Exact searches: $p = 1$ for $m \le 20{,}000$; $p = 2, 3$ to $n = 20{,}000$; no solutions & \texttt{period1\_cycles.py}, \texttt{period2\_cycles.py}, \texttt{period3\_cycles.py}; \texttt{cycles.md} \S12.2.3, \S12.5.3, \S12.7.5 \\
Theorem~\ref{thm:staircase}(i) & \texttt{staircase\_gamma\_upper.py}: $45$ checks, $0$ failures; \texttt{staircase\_allp\_\allowbreak construction.py} and \texttt{staircase\_allp\_\allowbreak diophantine.py}: the sweep, the bracket, the window counts, and the construction at every period from $3$ to $32$ apart from $p = 4$, including $3 \le p \le 15$; exact integer arithmetic at every size decision & \texttt{staircase\_gamma\_upper.py}; \texttt{staircase\_allp\_\allowbreak construction.py}, \texttt{staircase\_allp\_\allowbreak diophantine.py}; \texttt{cycles.md} \S12.8.6 \\
Theorem~\ref{thm:staircase}(i), $p \in \{2,4\}$ & The exhibitions & \texttt{staircase\_allp.py}; \texttt{cycles.md} \S12.8.6.4 \\
Theorem~\ref{thm:staircase}(ii)--(iii) & $188$ checks, $0$ failures: the sandwich and the trim on $30{,}000$ random configurations, the two-distance structure, $p_1(X_0)$ for $X_0 = 2^j$, $j \le 200$, the construction end to end at $95$ pairs $(X_0,p)$; and an independent re-implementation, $43$ checks, $0$ failures & \texttt{staircase\_x0\_threshold.py}; \texttt{x0\_threshold\_\allowbreak independent\_check.py}; \texttt{cycles.md} \S12.8.7 \\
The values of $p_1(X_0)$ (``The threshold, explicitly'') & The values at $X_0 = 2^j$, at $5\cdot 2^{60}$ and at $2^{71}$; the ranges follow by monotonicity & \texttt{staircase\_x0\_threshold.py}; \texttt{cycles.md} \S12.8.7.4 \\
Remark~\ref{rem:crash} & The table of the unit-block condition at $p_1(X_0)$ and $p_1(X_0) + 20$, $2^{20} \le X_0 \le 2^{200}$ & output of \texttt{staircase\_x0\_threshold.py}; \texttt{cycles.md} Remark~12.8.7.5 \\
Theorem~\ref{thm:staircase}(iv) & The $p = 7$ staircase; the seeded random search at $p = 6$ ($84$ size-passers, each clearing the period-$3$ constant by at least $0.698$) & \texttt{uniform\_trim.py}; \texttt{cycles.md} Remark~12.8.3 \\
Calibration (Appendix~\ref{app:aeh}) & The eight-round campaign; the fourth anomaly & \texttt{aeh\_calibration.py}; \texttt{aeh\_anomaly.py}; \texttt{aeh.md} \S13.4, \S13.5 \\
\hline
\end{longtable}
}

\paragraph{Statements proved in the record and not printed here.} The paper names the record section for each: the second display of Theorem~\ref{thm:uniform} (\texttt{cycles.md} Corollary~12.8.2, after \cite[Lemma~14, Theorem~3(d)]{sdw}); the details of Theorem~\ref{thm:smallp} (the \texttt{cycles.md} sections named in its outline); the comparisons of ``The threshold, explicitly'' (\texttt{cycles.md} Remark~12.8.7.6); the integer form of the elimination identity (\texttt{cycles.md} Remark~12.6.1.7); the deep-entry clause of the dictionary with $m$-cycles (\texttt{cycles.md} \S12.4); the depth renewal law (\texttt{aeh.md} \S13.6); the normalisation of $\mathrm{Syrac}$ (\texttt{aeh.md} \S13.6.5); the joint labelled law and the two-sided measure (\texttt{aeh.md} \S13.6.3, \S13.6.5); the equivalent form of Hypothesis~\ref{hyp:aeh} (\texttt{aeh.md} Theorem~13.6.4); the single-limit argument (\texttt{aeh.md} \S13.6.6); the base case and the exponent mean (\texttt{aeh.md} Lemma~13.2.4, with \S13.2.3); and the dyadic-scale density statement (\texttt{aeh.md} Proposition~13.2.5).

\section{Concordance with v3, and notation changes}\label{app:concordance}

v3 numbers its statements by section; v4 moves every formal statement into an appendix, in the same order, so that a v3 number becomes a v4 number by a change of the leading section number to a letter.

\begin{center}
\small
\begin{tabular}{lll@{\qquad}lll}
\hline
v3 & v4 & & v3 & v4 & \\
\hline
Def 2.1 & Def \ref{def:reduced} & reduced map & Prop 4.1 & Prop \ref{prop:elim} & elimination \\
Prop 2.2 & Prop \ref{prop:block} & blocks & Cor 4.2 & Cor \ref{cor:size} & size balance \\
Thm 2.3 & Thm \ref{thm:equiv} & equivalence & Lemma 4.3 & Lemma \ref{lem:ceiling} & ceiling \\
Def 3.1 & Def \ref{def:anchor} & anchor & Thm 4.4 & Thm \ref{thm:smallp} & periods 1--3 \\
Lemma 3.2 & Lemma \ref{lem:squaring} & squaring & Thm 4.5 & Thm \ref{thm:uniform} & uniform trim \\
Thm 3.3 & Thm \ref{thm:vlaw} & valuation law & Thm 4.6 & Thm \ref{thm:staircase} & staircase \\
Lemma 3.4 & Lemma \ref{lem:absorption} & absorption & Remark 4.7 & Remark \ref{rem:staircase} & divergence \\
Thm 3.5 & Thm \ref{thm:depth} & depth closure & --- & Remark \ref{rem:crash} & crash (new) \\
Remark 3.6 & Remark \ref{rem:verify1} & checks & Hyp 5.1 & Hyp \ref{hyp:aeh} & AEH \\
Thm 3.7 & Thm \ref{thm:deltaM} & increment & Lemma 5.2 & Lemma \ref{lem:routing} & routing \\
Thm 3.8 & Thm \ref{thm:onestep} & trichotomy & eq.\ (1) & eq.\ \eqref{eq:isometry} & isometry \\
Heur.\ 3.9 & Heur.\ \ref{prop:budget} & digit budget & --- & eqs.\ \eqref{eq:term}, \eqref{eq:element} & new \\
Remark 3.10 & Remark \ref{rem:verify2} & checks & & & \\
\hline
\end{tabular}
\end{center}

v3's sections map as follows: \S1 (introduction, related work and provenance) to Sections~\ref{sec:intro} and~\ref{sec:credit} and Appendix~\ref{app:related}; \S2 to Sections~\ref{sec:block}--\ref{sec:reduced} and Appendix~\ref{app:reduced}; \S3 to Section~\ref{sec:anchor} and Appendix~\ref{app:laws}; \S4 to Section~\ref{sec:cycles} and Appendix~\ref{app:cycles}; \S5 to Section~\ref{sec:aeh} and Appendix~\ref{app:aeh}; \S6 (discussion) to Section~\ref{sec:open}; Appendix~A to Appendix~\ref{app:method}.

\paragraph{Renamed symbols.} The symbols below are renamed; no meaning changed. ($\eta$ now names only the quantity $\phi(s_c)$ in the proof of Theorem~\ref{thm:staircase}.)
\begin{center}
\small
\begin{tabular}{>{\raggedright\arraybackslash}p{0.3\textwidth}>{\raggedright\arraybackslash}p{0.34\textwidth}>{\raggedright\arraybackslash}p{0.26\textwidth}}
\hline
v3 & v4 & meaning \\
\hline
$\varepsilon$ (Theorem 3.5) & $\xi$ (Theorem~\ref{thm:depth}) & the anchor displacement \\
$\varepsilon_t$ (Corollary 4.2, Lemma 4.3) & $\chi_t$ (Corollary~\ref{cor:size}, Lemma~\ref{lem:ceiling}) & the block ratio $(2^{s_t}-1)/(3^{d_t}\w_t)$ \\
$T_j$ (proof of Theorem 4.5) & $\mu_j$ (also in the proof of Theorem~\ref{thm:staircase} and Remark~\ref{rem:crash}) & partial sums of entry depths \\
$S$ (Theorem 4.6(iv)) & $\stot$ (Theorem~\ref{thm:staircase}(iv) and its proof) & $K - n$, the total exit valuation \\
$T_N$ (Hypothesis 5.1) & $H_N$ & the block horizon \\
$\Lambda_N$ (Hypothesis 5.1) & $E_N$ & the exponent budget \\
$H$ (\S5) & $\Ent$ & the binary entropy \\
$\eta$ (\S5) & $\zeta$ & a fixed positive slack \\
$(r,1)$ (Lemma 5.2) & $[r,1]$ & the class $\w \equiv r \pmod 8$, $d = 1$ \\
\hline
\end{tabular}
\end{center}
Unchanged: $T$ is the odd-to-odd map throughout; $\Lambda$ is the linear form $K\log 2 - n\log 3$; $\varepsilon$ is a tolerance; $M_t$ and $S_t$ (Proposition~\ref{prop:elim}) keep the record's notation.

\section{Symbol index}\label{app:symbols}

Every symbol of the paper, with its meaning and where it is defined. A symbol with two meanings is listed once per meaning; the context (the appendix) decides which is meant. Main-text symbols are also in ``Symbols at a glance'', before Section~\ref{sec:intro}.

{\small
\begin{longtable}{>{\raggedright\arraybackslash}p{0.17\textwidth}>{\raggedright\arraybackslash}p{0.53\textwidth}>{\raggedright\arraybackslash}p{0.22\textwidth}}
\hline
Symbol & Meaning & Defined \\
\hline
\endhead
\multicolumn{3}{l}{\emph{The map, blocks and states}} \\
$T$, $\mathrm{Syr}$ & the odd-to-odd Collatz map (Tao's Syracuse map) & App.~\ref{app:reduced} \\
$\vt{\cdot}$, $\vth{\cdot}$ & the number of times $2$, resp.\ $3$, divides a number & App.~\ref{app:reduced} \\
$x$ & an odd positive integer; in App.~\ref{app:aeh}, the sampled start & \S\ref{sec:intro}; Hyp.~\ref{hyp:aeh} \\
$x_j$ & the $j$-th iterate inside a block & Prop.~\ref{prop:block} \\
$m$ & entry depth: $\vt{x+1}$ & Prop.~\ref{prop:block} \\
$u$ & the odd part of $x+1$ & Prop.~\ref{prop:block} \\
$a$ & $\vth{u}$, the factors of $3$ in $u$ & Def.~\ref{def:reduced} \\
$\w$ & odd core, prime to $3$ & Def.~\ref{def:reduced} \\
$d$ & depth, $m + a$ & Def.~\ref{def:reduced} \\
$R$ & the projection $x \mapsto (\w, d)$ & Def.~\ref{def:reduced} \\
$F$ & the reduced map & Def.~\ref{def:reduced} \\
$A$ & $3^d\w - 1$ & Def.~\ref{def:reduced} \\
$s$ & exit valuation, $\vt{A}$ & Def.~\ref{def:reduced} \\
$x_{\mathrm{exit}}$ & the block's exit, $A/2^s$ & Def.~\ref{def:reduced} \\
$C$ & $A + 2^s$ & Def.~\ref{def:reduced} \\
$\sigma$ & $\vt{C} = s + m_+$; in \S\ref{sec:aeh} and App.~\ref{app:aeh} ``$5\sigma$'' also means five standard deviations & Def.~\ref{def:reduced} \\
$a_+$ & absorption, $\vth{C}$ & Def.~\ref{def:reduced} \\
$m_+$, $\w_+$, $d_+$ & the next entry depth, core and depth & Def.~\ref{def:reduced} \\
$e_t$, $S_t$ & the block entries and the states of an $F$-cycle (proof of Thm.~\ref{thm:equiv} only; not the $S_t$ of Prop.~\ref{prop:elim}) & proof of Thm.~\ref{thm:equiv} \\
\multicolumn{3}{l}{\emph{The exact laws (Appendix~\ref{app:laws})}} \\
$\Z_2$ & the $2$-adic integers & App.~\ref{app:laws} \\
$N(u)$ & for $u \equiv 1 \pmod 8$, the $t \in \Z_2$ with $9^t = u^{-1}$ & App.~\ref{app:laws} \\
$M(\w)$ & the anchor, $N(\w^2)$ & Def.~\ref{def:anchor} \\
$Y$ & $3^d\w$ & Lemma~\ref{lem:squaring} \\
$h(s)$ & $\vth{2^s - 1}$ & Lemma~\ref{lem:absorption} \\
$\xi$ & anchor displacement & Thm.~\ref{thm:depth} \\
$n$ & in Thm.~\ref{thm:depth}: half the depth ($d = 2n$ or $2n+1$); elsewhere see below & Thm.~\ref{thm:depth} \\
$\delta_s$ & $N(1 - 2^s)$ & Thm.~\ref{thm:depth} \\
$c$ & a target in $1 + 8\Z_2$ (target-shift lemma) & proof of Thm.~\ref{thm:depth} \\
$X$, $X'$ & $3^d\w$, $3^{d-1}\w$ (proof of Thm.~\ref{thm:depth}) & proof of Thm.~\ref{thm:depth} \\
$\Delta M$ & $M(\w_+) - M(\w)$ & Thm.~\ref{thm:deltaM} \\
$k$ & window depth; in App.~\ref{app:cycles}, also a convergent index & Thm.~\ref{thm:deltaM} \\
$q$, $r$ & exponents in facts (b), (c) of the proof of Thm.~\ref{thm:deltaM} (not the cycle's $q$) & proof of Thm.~\ref{thm:deltaM} \\
$s_+$ & the next exit valuation & Thm.~\ref{thm:onestep} \\
$W$ & a number of digits (Heuristic~\ref{prop:budget}) & Heur.~\ref{prop:budget} \\
\multicolumn{3}{l}{\emph{Cycles (Appendix~\ref{app:cycles})}} \\
$p$ & period: number of blocks & App.~\ref{app:cycles} \\
$t$, $r$ & block index in $\Z/p\Z$; rotation & App.~\ref{app:cycles} \\
$m_t$, $s_t$, $d_t$, $\w_t$, $a_t$ & the data of block $t$ & App.~\ref{app:cycles} \\
$n$ & $\sum_t m_t$ & App.~\ref{app:cycles} \\
$K$ & $\sum_t s_t + n$ & App.~\ref{app:cycles} \\
$q$ & $2^K - 3^n$ & App.~\ref{app:cycles} \\
$R_r$ & rotation numerator & Prop.~\ref{prop:elim} \\
$M_t$, $S_t$, $\sigma_j$ & partial sums of entry depths, of $\sigma$'s; $\sigma_j = s_j + m_{j+1}$ & Prop.~\ref{prop:elim} \\
$F_n$, $a$ & Tao's $n$-Syracuse offset; the profile's valuation word & after Prop.~\ref{prop:elim} \\
$\chi_t$ & block ratio $(2^{s_t}-1)/(3^{d_t}\w_t)$ & Cor.~\ref{cor:size} \\
$X$ & a lower bound on the cycle's elements & Cor.~\ref{cor:size} \\
$x_t$ & the exit of block $t$ & Cor.~\ref{cor:size} \\
$X_0$ & in Lemma~\ref{lem:ceiling}: the verified range; in Thm.~\ref{thm:staircase}: a size threshold & Lemma~\ref{lem:ceiling}; Thm.~\ref{thm:staircase} \\
$\LL$ & the base-$2$ logarithm of $3$, $1.58496\ldots$ & preamble \\
$\gamma$ & $K - \log_2 q$ & ``Configurations'' \\
$x_r$, $x_{\min}$ & block entries of a configuration; the least & ``Configurations'' \\
$\Lambda$ & $K\log 2 - n\log 3$ & ``Configurations'' \\
$\kp$ & the trim's exponent $(\LL-1)/((\LL)^p - 1)$ (the record's $\varepsilon_p$) & ``Configurations'' \\
$A$, $A_t$; $w(A)$, $m(A)$, $s(A)$ & arcs of consecutive blocks; arc weight, $m$-mass, $s$-mass & ``Configurations'' \\
$e_t$, $\Phi_r$, $w_j$, $r_0$ & exponents, recursion values, arc weights, reset index (proof of Thm.~\ref{thm:uniform}) & proof of Thm.~\ref{thm:uniform} \\
$\mu_j$ & partial sums of entry depths & proofs of Thm.~\ref{thm:uniform}, \ref{thm:staircase}; Rem.~\ref{rem:crash} \\
$n_0(p)$, $K_1(p)$ & the bound on $n$; Simons and de~Weger's constant & Thm.~\ref{thm:uniform} and after \\
$\phi$ & $-\log_2(1 - 2^{-x})$ & Thm.~\ref{thm:staircase} \\
$c_p$ & a trim constant & Thm.~\ref{thm:staircase}(i) \\
$\lambda_0$, $g$ & $\LL/(\LL-1) + 2^{-20}$; $1 + \max(\log_2(X_0+1), \lambda_0)$ & Thm.~\ref{thm:staircase}(iii) \\
$p_1(X_0)$ & the threshold period & proof of Thm.~\ref{thm:staircase}(iii) \\
$\rho_0$ & Rozier's quantity & Thm.~\ref{thm:staircase}(iii) \\
$\delta_n$ & $K - n\LL$ for $K = \lceil n\LL\rceil$ & proof of Thm.~\ref{thm:staircase} \\
$W_j$, $j^\ast$ & partial weights; their minimizer & proof of Thm.~\ref{thm:staircase}(ii) \\
$\stot$, $s_c$, $\eta$ & $K - n$; $\stot - p + 1$; $\phi(s_c)$ & proof of Thm.~\ref{thm:staircase} \\
$\Gamma'$, $\Gamma'_0$ & the threshold in (H1) & proof of Thm.~\ref{thm:staircase} \\
$f_r$, $b_r$, $X_r$ & the climb target, budget and comparison number & proof of Thm.~\ref{thm:staircase} \\
$\beta$ & in App.~\ref{app:cycles}: $1 + \eta$; in App.~\ref{app:aeh}: $2(2 - \LL)$, the per-block contraction rate & proofs of Thm.~\ref{thm:staircase}; App.~\ref{app:aeh} \\
$A_r$ & a real recursion & proof of Thm.~\ref{thm:staircase} \\
$\vartheta$, $\Gamma^\ast$, $\delta_{\mathrm{hi}}$, $\delta_{\mathrm{lo}}$ & the constants of part (i) & proof of Thm.~\ref{thm:staircase}(i) \\
$P_j/Q_j$, $\Theta_j$, $\Delta$, $N_\ast$, $n^\ast$ & convergents of $\LL$, their errors, the gap target, the window length, its end & proof of Thm.~\ref{thm:staircase}(iii) \\
$j_1$ & a unit block & Rem.~\ref{rem:crash} \\
\multicolumn{3}{l}{\emph{The hypothesis (Appendix~\ref{app:aeh})}} \\
$k$, $D$ & window depth; cap & ``The window law'' \\
$W_{k,D}$, $u \wedge D$ & the capped window; $\min(u, D)$ & ``The window law'' \\
$\hat B$, $B$ & the two-sided and one-sided Bernoulli laws on letters & ``The window law''; Hyp.~\ref{hyp:aeh} \\
$\pi_{k,D}$, $\pi^{(L)}_{k,D}$ & the window law; its $L$-block law & ``The window law''; after Hyp.~\ref{hyp:aeh} \\
$\mathrm{Syrac}$ & Tao's Syracuse random variable & ``The window law'' \\
$\tau$, $\theta$ & budget rate; block horizon rate & Hyp.~\ref{hyp:aeh} \\
$N$, $b$ & the sampling scale; $\lfloor \log_2 N\rfloor$ & Hyp.~\ref{hyp:aeh} \\
$E_N$, $H_N$ & exponent budget; block horizon & Hyp.~\ref{hyp:aeh} \\
$\ell_n$, $\tilde\ell_n$, $\dagger$ & the letter at block $n$; the tallied symbol; the blank & Hyp.~\ref{hyp:aeh} \\
$S_n$ & exponent spent before block $n$ & Hyp.~\ref{hyp:aeh} \\
$w$, $L$, $f_N(w,x)$, $B[w]$ & a word, its length, its frequency, its model value & Hyp.~\ref{hyp:aeh} \\
$\varepsilon$ & a tolerance (also in Inselmann's envelope, App.~\ref{app:related}) & Hyp.~\ref{hyp:aeh} \\
$X_N$ & a level with $\log X_N = o(\log N)$ & ``The budget and admissibility'' \\
$\mathbb{E}_B$ & mean under $B$ & ``The budget and admissibility'' \\
$J$, $S$, $\zeta$ & window exponent; total exponent of an itinerary; a fixed slack & ``The unconditional base case'' \\
$\Ent$ & binary entropy & ``The unconditional base case'' \\
$[r,1]$ & the class $\w \equiv r \pmod 8$, $d = 1$ & Lemma~\ref{lem:routing} \\
$\nu$ & a per-start empirical law & ``Calibration'' \\
\hline
\end{longtable}
}

\begin{thebibliography}{99}
\small
\bibitem{terras} Riho Terras, \emph{A stopping time problem on the positive integers}, Acta Arith.\ 30 (1976), no.~3, 241--252; DOI 10.4064/aa-30-3-241-252.
\bibitem{everett} C.~J.~Everett, \emph{Iteration of the number-theoretic function
  $f(2n) = n$, $f(2n+1) = 3n+2$}, Adv.\ Math.\ 25 (1977), 42--45.
\bibitem{korec} I.~Korec, \emph{A density estimate for the $3x+1$ problem},
  Math.\ Slovaca 44 (1994), no.\ 1, 85--89.
\bibitem{tao} T.~Tao, \emph{Almost all orbits of the Collatz map attain almost
  bounded values}, Forum Math.\ Pi 10 (2022), Paper No.\ e12, 56 pp.;
  arXiv:1909.03562.
\bibitem{inselmann} M.~Inselmann, \emph{An approximation of the Collatz map and a
  lower bound for the average total stopping time}, arXiv:2402.03276 (2024).
\bibitem{wirsching} G.~J.~Wirsching, \emph{The Dynamical System Generated by the
  $3n+1$ Function}, Lecture Notes in Math.\ 1681, Springer-Verlag, Berlin, 1998.
\bibitem{thomas} A.~Thomas, \emph{A non-uniform distribution property of most
  orbits, in case the $3x+1$ conjecture is true}, Acta Arith.\ \textbf{178}
  (2017), no.~2, 125--134; arXiv:1512.05852.
\bibitem{steiner} R.~P.~Steiner, \emph{A theorem on the Syracuse problem}, Proc.\ 7th Manitoba Conf.\ Numerical Math.\ and Computing (1977), Congr.\ Numer.\ 20, Utilitas Math., Winnipeg, 1978, pp.\ 553--559.
\bibitem{rozier90} O.~Rozier, \emph{D\'emonstration de l'absence de cycles d'une certaine forme pour le probl\`eme de Syracuse}, Singularit\'e 1 (1990), no.~3, 9--10.
\bibitem{rozier91} O.~Rozier, \emph{Probl\`eme de Syracuse: ``majorations'' \'el\'ementaires des cycles}, Singularit\'e 2 (1991), no.~5, 8--11.
\bibitem{bohm} C.~B\"ohm, G.~Sontacchi, \emph{On the existence of cycles of given length in integer sequences like $x_{n+1} = x_n/2$ if $x_n$ even, and $x_{n+1} = 3x_n + 1$ otherwise}, Atti Accad.\ Naz.\ Lincei Rend.\ Cl.\ Sci.\ Fis.\ Mat.\ Natur.\ (8) 64 (1978), no.~3, 260--264.
\bibitem{hh} L.~Halbeisen, N.~Hungerb\"uhler, \emph{Optimal bounds for the length of rational Collatz cycles}, Acta Arith.\ 78 (1997), no.~3, 227--239; DOI 10.4064/aa-78-3-227-239.
\bibitem{brox} T.~Brox, \emph{Collatz cycles with few descents}, Acta Arith.\ 92 (2000), no.~2, 181--188; DOI 10.4064/aa-92-2-181-188.
\bibitem{ladue} M.~D.~LaDue, \emph{Clusters of integers with equal total stopping times in the $3x+1$ problem}, Fibonacci Quart.\ 56 (2018), no.~2, 156--162; DOI 10.1080/00150517.2018.12427710; arXiv:1709.02979.
\bibitem{williams} J.~Williams, \emph{A Coordinate System for Collatz Dynamics}, arXiv:2607.01718v1 (2 July 2026); DOI 10.48550/arXiv.2607.01718.
\bibitem{sos} V.~T.~S\'os, \emph{On the distribution mod~1 of the sequence $n\alpha$}, Ann.\ Univ.\ Sci.\ Budapest.\ E\"otv\"os Sect.\ Math.\ 1 (1958), 127--134.
\bibitem{swierczkowski} S.~\'Swierczkowski, \emph{On successive settings of an arc on the circumference of a circle}, Fund.\ Math.\ 46 (1958), 187--189; DOI 10.4064/fm-46-2-187-189.
\bibitem{eliahou} S.~Eliahou, \emph{The 3x+1 problem: new lower bounds on nontrivial cycle lengths}, Discrete Math.\ 118 (1993), 45--56.
\bibitem{sdw} J.~Simons, B.~de~Weger, \emph{Theoretical and computational bounds for $m$-cycles of the $3n+1$ problem}, Acta Arith.\ 117 (2005), no.~1, 51--70; DOI 10.4064/aa117-1-3.
\bibitem{hercher} C.~Hercher, \emph{There are no Collatz $m$-cycles with $m \le 91$}, J.\ Integer Seq.\ 26 (2023), Article 23.3.5; arXiv:2201.00406; corrigendum (14 June 2026), repairing a derivation in the proof of Theorem~21 with no change to any statement.
\bibitem{lagarias} J.~C.~Lagarias, \emph{The 3x+1 problem: an annotated bibliography}, I--II, arXiv math/0309224, math/0608208.
\bibitem{yu} K.~Yu, \emph{Linear forms in $p$-adic logarithms}, Acta Arith.\ 53 (1989), 107--186; \emph{Linear forms in $p$-adic logarithms II}, Compositio Math.\ 74 (1990), 15--113; \emph{Linear forms in $p$-adic logarithms III}, Compositio Math.\ 91 (1994), 241--276.
\bibitem{bl} Y.~Bugeaud, M.~Laurent, \emph{Minoration effective de la distance $p$-adique entre puissances de nombres alg\'ebriques}, J.\ Number Theory 61 (1996), 311--342.
\bibitem{rhin} G.~Rhin, \emph{Approximants de Pad\'e et mesures effectives d'irrationalit\'e}, in: S\'eminaire de Th\'eorie des Nombres, Paris 1985--86, Progr.\ Math.\ 71, Birkh\"auser, Boston, 1987, pp.\ 155--164 (the Proposition used here is on p.~160); see also Q.~Wu, \emph{On the linear independence measure of logarithms of rational numbers}, Math.\ Comp.\ 72 (2003), no.\ 242, 901--911.
\bibitem{barina} D.~Barina, \emph{Improved verification limit for the convergence of the Collatz conjecture}, J.\ Supercomputing 81 (2025), no.~7, article 810; DOI 10.1007/s11227-025-07337-0.
\bibitem{merle} E.~Merle, \emph{On the non-existence of non-trivial Collatz cycles: a conditional formal proof in Lean 4 with documented structural obstructions}, preprint, Zenodo (2026), the version of 26 April 2026 cited here at its version-specific DOI 10.5281/zenodo.19790406 (concept DOI 10.5281/zenodo.19790405).
\bibitem{rackl} D.~Rackl, \emph{Cycles in the Collatz sequence: an approach using $2$-adic arithmetic}, Bachelor's thesis, University of Klagenfurt, Department of Mathematics, 22 September 2021; \url{https://benjamin-hackl.at/downloads/theses/2021-bachelor-rackl.pdf}.
\bibitem{llmcollatz} E.~Y.~Chang, \emph{Exploring Collatz Dynamics with Human--LLM Collaboration}, arXiv:2603.11066 (2026).
\end{thebibliography}

\end{document}
